\documentclass[10pt]{extarticle}
\usepackage[a4paper,margin=1.5cm]{geometry}
\usepackage[T1]{fontenc}
\usepackage{amsmath,amssymb}
\usepackage{graphicx}
\usepackage[export]{adjustbox}
\usepackage{multicol}
\usepackage{enumitem}
\usepackage{booktabs,array,tabularx}
\setlength{\columnsep}{0.7cm}
\setlength{\parindent}{0pt}
\setlength{\parskip}{3pt}
\pagestyle{empty}
\makeatletter
% figure registry (figdefs.tex is generated by the extract step)
\newcommand{\qfigdef}[4]{\@namedef{qf@file@#1}{#2}\@namedef{qf@w@#1}{#3}\@namedef{qf@wide@#1}{#4}}
\newcommand{\qf@check}[1]{\@ifundefined{qf@file@#1}{\PackageError{score2}{Unknown figure id #1}{}}{}}
\newcommand{\qf@img}[4][]{\includegraphics[width=#2\dimexpr\@nameuse{qf@w@#4}pt\relax,max width=#3,max height=0.45\textheight,#1]{\@nameuse{qf@file@#4}}}
% Full-width questions: finalize wraps every question that contains a wide figure in \qfullbegin...\qfullend,
% so the WHOLE question (stem, figures, options) is set across the page in reading order, then the two columns
% resume. The question's head - its number, stem and first figure (\qfullsplit marks the end of the head) -
% is kept on one page, so a number is never left at a page foot with its figure overleaf; the rest may break
% across pages normally, so no big gaps are left. Without a \qfullsplit, the whole question is kept together
% if it is under 45% of a page. Never write these yourself.
\newif\ifqf@full
\newif\ifqf@boxed
\newsavebox{\qf@box}
\newcommand{\qf@open}{\global\qf@fulltrue\par\global\qf@boxedtrue
  \setbox\qf@box\vbox\bgroup\hsize\textwidth\linewidth\textwidth}
\newcommand{\qf@emit}[1]{\par\egroup\global\qf@boxedfalse
  \ifdim\dimexpr\ht\qf@box+\dp\qf@box\relax<#1\textheight\noindent\usebox{\qf@box}\par\else\unvbox\qf@box\fi}
\newcommand{\qfullsplit}{\ifqf@boxed\qf@emit{0.85}\fi}
\newcommand{\qf@close}{\ifqf@boxed\qf@emit{0.45}\fi\global\qf@fullfalse}
\newcommand{\qfullbegin}{\end{multicols}\qf@open}
\newcommand{\qfullend}{\qf@close\begin{multicols}{2}}
\newcommand{\qfullstart}{\qf@open}
\newcommand{\qfulldone}{\qf@close}
% \qfig[scale]{id}: block figure, centred, directly where it is written (size set by the tool).
% Fallback if a wide figure is not inside a full-width question: break out of the columns in place - no floats.
\newcommand{\qfig}[2][1]{\qf@check{#2}%
  \ifqf@full\qf@blk{#1}{#2}\else
  \ifnum\@nameuse{qf@wide@#2}=1 \expandafter\@firstoftwo\else\expandafter\@secondoftwo\fi
  {\end{multicols}\begin{center}\qf@img{#1}{\textwidth}{#2}\end{center}\begin{multicols}{2}}%
  {\qf@blk{#1}{#2}}\fi}
\newcommand{\qf@blk}[2]{\par\smallskip{\centering\qf@img{#1}{\linewidth}{#2}\par}\smallskip}
% \qfigi[scale]{id}: inline figure (inside option lists / tables / minipages); never breaks columns.
\newcommand{\qfigi}[2][1]{\qf@check{#2}\qf@img[valign=c]{#1}{\linewidth}{#2}}
\makeatother
% structure macros - lint reads these to check completeness and figure placement
\newcommand{\papersec}[1]{\par\bigskip{\centering\Large\bfseries PAPER-#1\par}\medskip}
\newcommand{\partsec}[2]{\par\medskip\noindent{\large\bfseries PART-#1: #2}\par\smallskip}
\newcommand{\sect}[1]{\par\smallskip\noindent\textbf{#1}\par}
\newcommand{\qn}[1]{\par\smallskip\noindent\textbf{#1.}~}
\newcommand{\paraq}[2]{\par\smallskip\noindent\textbf{Paragraph for Questions #1 and #2}\par}
\newenvironment{qoptions}{\begin{itemize}[leftmargin=1.8em,itemsep=1pt,topsep=2pt,label={}]}{\end{itemize}}
\newcommand{\opt}[2]{\item[(#1)] #2}
% figure registry (generated by score2_tools.py finalize)
\qfigdef{p1_phy_s1_q01_1}{Score Advanced - Test 28 (27 Apr 2024)-img/p1_phy_s1_q01_1.png}{161.1}{0}
\qfigdef{p1_phy_s1_q03_1}{Score Advanced - Test 28 (27 Apr 2024)-img/p1_phy_s1_q03_1.png}{108.0}{0}
\qfigdef{p1_phy_s1_q04_1}{Score Advanced - Test 28 (27 Apr 2024)-img/p1_phy_s1_q04_1.png}{133.6}{0}
\qfigdef{p1_phy_s1_q05_1}{Score Advanced - Test 28 (27 Apr 2024)-img/p1_phy_s1_q05_1.png}{120.4}{0}
\qfigdef{p1_phy_s2_q06_1}{Score Advanced - Test 28 (27 Apr 2024)-img/p1_phy_s2_q06_1.png}{142.3}{0}
\qfigdef{p1_phy_s2_q07_1}{Score Advanced - Test 28 (27 Apr 2024)-img/p1_phy_s2_q07_1.png}{160.7}{0}
\qfigdef{p1_phy_s2_q09_1}{Score Advanced - Test 28 (27 Apr 2024)-img/p1_phy_s2_q09_1.png}{170.7}{0}
\qfigdef{p1_phy_s3_q01_1}{Score Advanced - Test 28 (27 Apr 2024)-img/p1_phy_s3_q01_1.png}{87.0}{0}
\qfigdef{p1_phy_s3_q02_1}{Score Advanced - Test 28 (27 Apr 2024)-img/p1_phy_s3_q02_1.png}{55.6}{0}
\qfigdef{p1_phy_s3_q04_1}{Score Advanced - Test 28 (27 Apr 2024)-img/p1_phy_s3_q04_1.png}{74.0}{0}
\qfigdef{p1_phy_s3_q04_2}{Score Advanced - Test 28 (27 Apr 2024)-img/p1_phy_s3_q04_2.png}{76.8}{0}
\qfigdef{p1_phy_s3_q05_1}{Score Advanced - Test 28 (27 Apr 2024)-img/p1_phy_s3_q05_1.png}{103.1}{0}
\qfigdef{p1_phy_s3_q07_1}{Score Advanced - Test 28 (27 Apr 2024)-img/p1_phy_s3_q07_1.png}{103.6}{0}
\qfigdef{p1_phy_s3_q08_1}{Score Advanced - Test 28 (27 Apr 2024)-img/p1_phy_s3_q08_1.png}{304.2}{1}
\qfigdef{p1_chem_s1_q01_1}{Score Advanced - Test 28 (27 Apr 2024)-img/p1_chem_s1_q01_1.png}{192.0}{0}
\qfigdef{p1_chem_s1_q01_2}{Score Advanced - Test 28 (27 Apr 2024)-img/p1_chem_s1_q01_2.png}{182.0}{0}
\qfigdef{p1_chem_s1_q01_3}{Score Advanced - Test 28 (27 Apr 2024)-img/p1_chem_s1_q01_3.png}{205.7}{0}
\qfigdef{p1_chem_s1_q01_4}{Score Advanced - Test 28 (27 Apr 2024)-img/p1_chem_s1_q01_4.png}{245.2}{0}
\qfigdef{p1_chem_s1_q01_5}{Score Advanced - Test 28 (27 Apr 2024)-img/p1_chem_s1_q01_5.png}{159.4}{0}
\qfigdef{p1_chem_s1_q02_1}{Score Advanced - Test 28 (27 Apr 2024)-img/p1_chem_s1_q02_1.png}{230.7}{0}
\qfigdef{p1_chem_s1_q02_2}{Score Advanced - Test 28 (27 Apr 2024)-img/p1_chem_s1_q02_2.png}{28.6}{0}
\qfigdef{p1_chem_s1_q02_3}{Score Advanced - Test 28 (27 Apr 2024)-img/p1_chem_s1_q02_3.png}{33.3}{0}
\qfigdef{p1_chem_s2_q07_1}{Score Advanced - Test 28 (27 Apr 2024)-img/p1_chem_s2_q07_1.png}{245.2}{0}
\qfigdef{p1_chem_s3_q01_1}{Score Advanced - Test 28 (27 Apr 2024)-img/p1_chem_s3_q01_1.png}{92.5}{0}
\qfigdef{p1_chem_s3_q02_1}{Score Advanced - Test 28 (27 Apr 2024)-img/p1_chem_s3_q02_1.png}{182.2}{0}
\qfigdef{p2_phy_s1_q01_1}{Score Advanced - Test 28 (27 Apr 2024)-img/p2_phy_s1_q01_1.png}{97.9}{0}
\qfigdef{p2_phy_s1_q02_1}{Score Advanced - Test 28 (27 Apr 2024)-img/p2_phy_s1_q02_1.png}{150.6}{0}
\qfigdef{p2_phy_s1_q03_1}{Score Advanced - Test 28 (27 Apr 2024)-img/p2_phy_s1_q03_1.png}{150.6}{0}
\qfigdef{p2_phy_s1_q04_1}{Score Advanced - Test 28 (27 Apr 2024)-img/p2_phy_s1_q04_1.png}{117.1}{0}
\qfigdef{p2_phy_s1_q05_1}{Score Advanced - Test 28 (27 Apr 2024)-img/p2_phy_s1_q05_1.png}{171.7}{0}
\qfigdef{p2_phy_s1_q06_1}{Score Advanced - Test 28 (27 Apr 2024)-img/p2_phy_s1_q06_1.png}{99.4}{0}
\qfigdef{p2_phy_s2_q07_1}{Score Advanced - Test 28 (27 Apr 2024)-img/p2_phy_s2_q07_1.png}{101.8}{0}
\qfigdef{p2_phy_s2_q08_1}{Score Advanced - Test 28 (27 Apr 2024)-img/p2_phy_s2_q08_1.png}{95.4}{0}
\qfigdef{p2_phy_s2_q08_2}{Score Advanced - Test 28 (27 Apr 2024)-img/p2_phy_s2_q08_2.png}{70.1}{0}
\qfigdef{p2_phy_s2_q08_3}{Score Advanced - Test 28 (27 Apr 2024)-img/p2_phy_s2_q08_3.png}{87.4}{0}
\qfigdef{p2_phy_s2_q08_4}{Score Advanced - Test 28 (27 Apr 2024)-img/p2_phy_s2_q08_4.png}{70.8}{0}
\qfigdef{p2_phy_s2_q09_1}{Score Advanced - Test 28 (27 Apr 2024)-img/p2_phy_s2_q09_1.png}{158.7}{0}
\qfigdef{p2_phy_s2_q10_1}{Score Advanced - Test 28 (27 Apr 2024)-img/p2_phy_s2_q10_1.png}{75.4}{0}
\qfigdef{p2_phy_s2_q10_5}{Score Advanced - Test 28 (27 Apr 2024)-img/p2_phy_s2_q10_5.png}{108.7}{0}
\qfigdef{p2_phy_s2_q10_2}{Score Advanced - Test 28 (27 Apr 2024)-img/p2_phy_s2_q10_2.png}{73.2}{0}
\qfigdef{p2_phy_s2_q10_6}{Score Advanced - Test 28 (27 Apr 2024)-img/p2_phy_s2_q10_6.png}{105.3}{0}
\qfigdef{p2_phy_s2_q10_3}{Score Advanced - Test 28 (27 Apr 2024)-img/p2_phy_s2_q10_3.png}{74.2}{0}
\qfigdef{p2_phy_s2_q10_7}{Score Advanced - Test 28 (27 Apr 2024)-img/p2_phy_s2_q10_7.png}{100.1}{0}
\qfigdef{p2_phy_s2_q10_4}{Score Advanced - Test 28 (27 Apr 2024)-img/p2_phy_s2_q10_4.png}{77.3}{0}
\qfigdef{p2_phy_s2_q10_8}{Score Advanced - Test 28 (27 Apr 2024)-img/p2_phy_s2_q10_8.png}{95.8}{0}
\qfigdef{p2_phy_s2_q11_1}{Score Advanced - Test 28 (27 Apr 2024)-img/p2_phy_s2_q11_1.png}{118.8}{0}
\qfigdef{p2_phy_s2_q12_1}{Score Advanced - Test 28 (27 Apr 2024)-img/p2_phy_s2_q12_1.png}{109.2}{0}
\qfigdef{p2_phy_s3_q02_1}{Score Advanced - Test 28 (27 Apr 2024)-img/p2_phy_s3_q02_1.png}{90.1}{0}
\qfigdef{p2_phy_s3_q04_1}{Score Advanced - Test 28 (27 Apr 2024)-img/p2_phy_s3_q04_1.png}{96.3}{0}
\qfigdef{p2_phy_s3_q05_1}{Score Advanced - Test 28 (27 Apr 2024)-img/p2_phy_s3_q05_1.png}{132.3}{0}
\qfigdef{p2_phy_s3_q05_2}{Score Advanced - Test 28 (27 Apr 2024)-img/p2_phy_s3_q05_2.png}{46.6}{0}
\qfigdef{p2_phy_s3_q05_3}{Score Advanced - Test 28 (27 Apr 2024)-img/p2_phy_s3_q05_3.png}{47.3}{0}
\qfigdef{p2_phy_s3_q05_4}{Score Advanced - Test 28 (27 Apr 2024)-img/p2_phy_s3_q05_4.png}{47.3}{0}
\qfigdef{p2_phy_s3_q05_5}{Score Advanced - Test 28 (27 Apr 2024)-img/p2_phy_s3_q05_5.png}{46.1}{0}
\qfigdef{p2_phy_s3_q06_1}{Score Advanced - Test 28 (27 Apr 2024)-img/p2_phy_s3_q06_1.png}{133.0}{0}
\qfigdef{p2_chem_s1_q01_1}{Score Advanced - Test 28 (27 Apr 2024)-img/p2_chem_s1_q01_1.png}{284.1}{1}
\qfigdef{p2_chem_s1_q02_A}{Score Advanced - Test 28 (27 Apr 2024)-img/p2_chem_s1_q02_A.png}{97.5}{0}
\qfigdef{p2_chem_s1_q02_B}{Score Advanced - Test 28 (27 Apr 2024)-img/p2_chem_s1_q02_B.png}{129.4}{0}
\qfigdef{p2_chem_s1_q02_C}{Score Advanced - Test 28 (27 Apr 2024)-img/p2_chem_s1_q02_C.png}{134.5}{0}
\qfigdef{p2_chem_s1_q02_D}{Score Advanced - Test 28 (27 Apr 2024)-img/p2_chem_s1_q02_D.png}{92.0}{0}
\qfigdef{p2_chem_s2_q07_1}{Score Advanced - Test 28 (27 Apr 2024)-img/p2_chem_s2_q07_1.png}{109.6}{0}
\qfigdef{p2_chem_s2_q07_2}{Score Advanced - Test 28 (27 Apr 2024)-img/p2_chem_s2_q07_2.png}{106.3}{0}
\qfigdef{p2_chem_s2_q07_3}{Score Advanced - Test 28 (27 Apr 2024)-img/p2_chem_s2_q07_3.png}{106.1}{0}
\qfigdef{p2_chem_s2_q07_4}{Score Advanced - Test 28 (27 Apr 2024)-img/p2_chem_s2_q07_4.png}{111.1}{0}
\qfigdef{p2_chem_s2_q08_1}{Score Advanced - Test 28 (27 Apr 2024)-img/p2_chem_s2_q08_1.png}{113.7}{0}
\qfigdef{p2_chem_s2_q08_2}{Score Advanced - Test 28 (27 Apr 2024)-img/p2_chem_s2_q08_2.png}{96.3}{0}
\qfigdef{p2_chem_s2_q08_3}{Score Advanced - Test 28 (27 Apr 2024)-img/p2_chem_s2_q08_3.png}{106.8}{0}
\qfigdef{p2_chem_s2_q08_4}{Score Advanced - Test 28 (27 Apr 2024)-img/p2_chem_s2_q08_4.png}{150.8}{0}
\qfigdef{p2_chem_s2_q08_5}{Score Advanced - Test 28 (27 Apr 2024)-img/p2_chem_s2_q08_5.png}{113.7}{0}
\qfigdef{p2_chem_s3_pre_1}{Score Advanced - Test 28 (27 Apr 2024)-img/p2_chem_s3_pre_1.png}{126.8}{0}
\qfigdef{p2_chem_s3_pre_2}{Score Advanced - Test 28 (27 Apr 2024)-img/p2_chem_s3_pre_2.png}{131.8}{0}
\qfigdef{p2_chem_s3_pre_3}{Score Advanced - Test 28 (27 Apr 2024)-img/p2_chem_s3_pre_3.png}{81.1}{0}
\qfigdef{p2_chem_s3_pre_4}{Score Advanced - Test 28 (27 Apr 2024)-img/p2_chem_s3_pre_4.png}{53.4}{0}
\qfigdef{p2_chem_s3_q01_1}{Score Advanced - Test 28 (27 Apr 2024)-img/p2_chem_s3_q01_1.png}{116.8}{0}
\qfigdef{p2_chem_s3_q01_4}{Score Advanced - Test 28 (27 Apr 2024)-img/p2_chem_s3_q01_4.png}{176.2}{0}
\qfigdef{p2_chem_s3_q01_2}{Score Advanced - Test 28 (27 Apr 2024)-img/p2_chem_s3_q01_2.png}{154.8}{0}
\qfigdef{p2_chem_s3_q01_3}{Score Advanced - Test 28 (27 Apr 2024)-img/p2_chem_s3_q01_3.png}{157.5}{0}
\raggedcolumns
\makeatletter\@beginparpenalty=10000\makeatother
\newcolumntype{L}[1]{>{\raggedright\arraybackslash}p{#1}}
\begin{document}
\begin{center}{\Large\bfseries Score Advanced --- Test 28 (27 Apr 2024)}\end{center}
\begin{multicols}{2}
\raggedcolumns
\papersec{1}
\partsec{1}{PHYSICS}
\sect{SECTION-II (i)}
\qn{1} A source of current i$_0\sin\omega$t, with i$_0$ as constant, is connected to the circuit shown in the diagram. The frequency $\omega$ is controllable.\par
\qfig{p1_phy_s1_q01_1}
The inductors L$_1$, L$_2$ and capacitor C$_1$, C$_2$ are all ideal. An ideal voltmeter is connected between points A and B. For very small $\omega$ (but non zero), reading of voltmeter is $\dfrac{\mathrm i_0\omega\mathrm L_1}{\mathrm x}$. Find x.
\qn{2} The rubber shell of a balloon has mass m = 585 $\dfrac\pi2$ gm. We fill methane (CH$_4$) in it at room temperature. The molar mass of air can be assumed to be 29 gm/mol. The balloon always remains spherical in shape and the pressure difference between inside \& outside due to its elastic nature can be neglected. At what diameter (in cm) will the balloon float in air? Thickness of balloon is negligible.\par
(p$_\mathrm{atm}$ = 10$^5$ Pa, R = $\dfrac{25}3$ J/mol K, T$_0$ = 27$^\circ$C.)
\qn{3} A parallel beam of light falls on a quarter cylinder of radius 9 cm as shown in figure (A). Refractive index of the material of the cylinder is $\sqrt3$. Maximum value of OP in meter shown in figure (B), so that rays don't suffer TIR at the curved surface is (consider refractive index of surrounding to be 1).\par
\qfig{p1_phy_s1_q03_1}
\qn{4} A monochromatic light of wavelength $\lambda_0$ and intensity I$_0$ is incident on YDSE experiment. The distance between slits S$_1$S$_2$ is `d' (d$\ll$D). The angle made by parallel beam light with the normal to the plane of slits is $\sin^{-1}\left(\dfrac{3\lambda}{2\mathrm d}\right)$. Now, if a transparent glass slab of thickness $\dfrac{3\lambda}{2(\mu-1)}$ is kept infront of one of the slits, where `$\mu$' is refractive index of glass slab, find the intensity (in W/m$^2$) at the geometrical centre `O' of the screen. (Given, d = 1mm, D = 2m, $\mu$ = 4/3, $\lambda$ = 4000 \AA\ and I$_0$ = 4.16 W/m$^2$)\par
\qfig{p1_phy_s1_q04_1}
\qn{5} What can be the maximum amplitude (in m) of the system so that there is no slipping between any of the blocks ? (Take g = 10 ms$^{-2}$).\par
\qfig{p1_phy_s1_q05_1}
\sect{SECTION-II (ii)}
\qn{6} A source of sound frequency 1000 Hz moves to the right with a speed of 50 m/s relative to the ground. To its right is a reflecting surface moving to the left with a speed of 20 m/s. There is a stationary observer to the left of the source as shown in figure. Speed of sound in air is 330 m/s. Then wavelength of reflected wave (in mm) is :-\par
\qfig{p1_phy_s2_q06_1}
\qn{7} Figure shows inverted T shaped fixed pipe. Liquid is released from rest from the position shown. Liquid is non-viscous. Velocity of fluid particle in vertical part of tube when level of liquid in it remains H/2 will be $\sqrt{\dfrac{\mathrm{xgH}}8}$. Find x.\par
\qfig{p1_phy_s2_q07_1}
\qn{8} A satellite weighing 3 tons flies in an approximately circular orbit near earth. The resistive force acting on it is 0.01 N. How much will its flight height (in km) decrease in a period of 10 day? (Round off to nearest integer) (Earth's radius = 6400 km, g = 10 m/s$^2$)
\qn{9} APQ is a rough surface. The portion AP is parabolic given by equation x$^2$ = 16y. The point `A' assumed to be origin. The portion PQ is a straight line of length 20$\sqrt3$ m, the coefficient of friction between block and surface (APQ) is $\dfrac1{\sqrt3}$. A block of mass 3 kg is moved on surface APQ by applying a force on it. The direction of force is always tangential to the surface APQ. The line PQ is tangent on the curve AP at P. If the block moves very slowly on the surface then find the work done by applying force in joule till it reaches Q.\par
\qfig{p1_phy_s2_q09_1}
\qn{10} A 2000 W resistance heater wire has thermal conductivity 15 W/mK, radius of 4 mm and length 0.5 m. It is used to boil water. The outer surface of wire has temperature of 105$^\circ$C. What is the temperature (in $^\circ$C) at the axis of the wire, if heat flows out only radially? Round off to nearest integer.
\sect{SECTION-III}
\qn{1} There is an insulating uniformly charged cylindrical shell of radius a lies with its axis along the z-axis. The shell's moment of inertia per unit length about the z-axis and the surface charge density are I and $\sigma$ respectively. The cylinder is placed in an external uniform magnetic field B$_{\mathrm{ex}}\hat{\mathrm z}$, and is initially at rest. Starting at t = 0 the external magnetic field is slowly reduced to zero. If the final angular velocity $\omega$ of the cylinder is $\omega=\dfrac{\sigma\pi\mathrm a^{\mathrm x}\mathrm B}{\mathrm I}\left(1+\dfrac{\mu_0\sigma^{\mathrm y}\mathrm a^{\mathrm z}\pi}{\mathrm I}\right)^{\mathrm d}$, find x + y + z + 2d.\par
\qfig{p1_phy_s3_q01_1}
\qn{2} We have a cube whose 3 adjacent faces are charged uniformly with surface charge density of $\sigma$. The remaining 3 faces are charged with a surface charge density of $-\sigma$. If the electric field at the centre of the cube is $\dfrac{\sigma}{\sqrt{\mathrm x}\epsilon_0}$, find x.\par
\qfig{p1_phy_s3_q02_1}
\qn{3} A sphere made of polymer is floating on the surface of water which is at a temperature of 20$^\circ$C and 1\% of the volume of the sphere is above the water. To what final temperature (in $^\circ$C) must the water be heated in order to submerge the polymer sphere completely ? Take coefficient of linear expansion of the polymer 2.30 $\times$ 10$^{-5}$/$^\circ$C, and coefficient of volume expansion of water 3.19 $\times$ 10$^{-4}$/$^\circ$C. If your answer is N, find value of $\dfrac{\mathrm N}{20}$.
\qn{4} We want to measure the resistance of a resistor. A voltmeter and ammeter are connected as shown in figure a \& b. Their readings are also indicated in figure. What is the voltage of the battery (in V)? The battery is ideal.\par
\qfigi{p1_phy_s3_q04_1}\hfill\qfigi{p1_phy_s3_q04_2}\par
\qn{5} A ball is thrown from the inclined plane (large) at an angle $\theta$ from inclined surface with an initial speed u$_0$. For the resulting projectile motion, after hitting the surface, ball rebounds at the same angle $\theta$, but with reduced speed $\dfrac{\mathrm u_0}{\beta}$. Its motion continues for a long time. If the magnitude of average velocity of the ball for entire duration of motion is u$_0$, then value of $\beta$ is :-\par
\qfig{p1_phy_s3_q05_1}
\qn{6} In the determination of Young's modulus by using Searle's method, a wire of length L = 2m and diameter d = 0.5 mm is used. For a load M = 2.5 kg, an extension $\ell$ = 0.25 mm in the length of the wire is observed. Quantities d and $\ell$ are measured using a screw gauge and a micrometer, respectively. They have the same pitch of 0.5 mm. The number of divisions on their circular scale is 100. The length of wire, acceleration due to gravity and mass M are known exactly. The percentage error in measurement of the Young's modulus is\underline{\hspace{2cm}}.
\qn{7} Two identical simple pendulums each of length $\ell$ are connected by a weightless spring as shown in figure. The force constant of the spring is k. In equilibrium, the pendulums are vertical and the spring is horizontal and underformed. If the time period of small oscillation of the linked pendulums, when they are deflected from their equilibrium positions through equal displacements in the same vertical plane in opposite direction and released is $2\pi\sqrt{\dfrac{\ell}{\left(\mathrm g+\frac{\mathrm{nk}\ell}{\mathrm m}\right)}}$, find n.\par
\qfig{p1_phy_s3_q07_1}
\qfullbegin
\qn{8} Six capacitors and four ideal batteries are connected in a circuit as shown in the figure. Initially all capacitors are uncharged and all switches were open. Now, all switches are closed simultaneously. Work done by $\varepsilon_1$ = 3V battery is\underline{\hspace{1cm}} (in $\mu$J) :\par
\qfig{p1_phy_s3_q08_1}
\qfullsplit
\raggedcolumns
\qfullend
\papersec{1}
\partsec{2}{CHEMISTRY}
\sect{SECTION-II (i)}
\qn{1} How many ether(s) will undergo cleavage on reaction with conc.HI ?\par
\qfig{p1_chem_s1_q01_1}
\qfig{p1_chem_s1_q01_2}
\qfig{p1_chem_s1_q01_3}
\qfig{p1_chem_s1_q01_4}
\qfig{p1_chem_s1_q01_5}
\qn{2} How many compounds are more reactive towards S$_\mathrm N$1 as well as S$_\mathrm N$2 than CH$_3-$CH$_2-$CH$_2-$Cl :\par
\qfig{p1_chem_s1_q02_1}
(iv) CH$_3-$CH$_2-\ddot{\mathrm N}$H$-$CH$_2-$Cl\quad (v) \qfigi[0.75]{p1_chem_s1_q02_2}\par
(vi) CH$_2$=CH$-$CH$_2-$Cl\quad (vii) \qfigi[0.75]{p1_chem_s1_q02_3}\par
\qn{3} At 300 K temperature and 8 atm pressure the compressibility factor for a real gas is 0.8. The volume (in litre) of 10 mole of the gas at same temperature and pressure is\par
[Given : R = 0.0821 L-atm mol$^{-1}$ K$^{-1}$]
\qn{4} The rate constant for the reaction N$_2$O$_4$(g) $\longrightarrow$ 2NO$_2$(g) is 92.12 $\times$ 10$^{-2}$ min$^{-1}$. The time (in seconds) required for the molar ratio of NO$_2$ (g) to N$_2$O$_4$ (g) in the reaction mixture to become 2 : 1 is\par
[log 2 = 0.3010, ln 10 = 2.303]
\qn{5} Sum of the spin only magnetic moment (in the units of B.M.) of the following complexes:\par
[Cu(NH$_3$)$_4$]$^{2+}$ ; [Fe(CN)$_6$]$^{3-}$ ; [NiCl$_4$]$^{2-}$ ; [CoF$_6$]$^{3-}$ ; [Fe(H$_2$O)$_5$(NO)]SO$_4$\par
[Given: $\sqrt2$ = 1.41 , $\sqrt3$ = 1.73 , $\sqrt5$ = 2.24 , $\sqrt7$ = 2.64 , $\sqrt{15}$ = 3.87 , $\sqrt{24}$ = 4.90]
\sect{SECTION-II (ii)}
\qn{6} How many are correctly matched for purification of organic compounds?\par
(a) Distillation is used for separation of chloroform \& aniline.\par
(b) Fractional distillation for crude oil in petroleum industry.\par
(c) Distillation under reduced pressure is used for liquid having very high boiling point.\par
(d) Distillation under reduced pressure is used for spent-eye in soap industry.\par
(e) Steam distillation for aniline water mixture.\par
(f) Paper chromatography is type of adsorption chromatography.\par
(g) Column chromatography is type of paper chromatography.\par
(h) Thin layer chromatography is type of adsorption chromatography.
\qn{7} How many of the following can be prepared by gabriel phthalimide synthesis in good yield.\par
\qfig{p1_chem_s2_q07_1}
\qn{8} An ionic compound is created with anions A \& B present at corners of cube in 3 : 1 ratio. All face$-$centers \& body centre contain cation Z. Atomic weight of A, B and Z are respectively 20, 40 \& 15 amu. Edge length of unit cell is 1 \AA, density = d(amu \AA$^{-3}$) and empirical formula of ionic compound is A$_\mathrm a$B$_\mathrm b$Z$_\mathrm c$, here a, b and c are smallest integer value. Then value of (a + b + c + d) considering unit cell is electrically neutral.
\qn{9} 1 mole of an ideal monoatomic gas initially at 1 atm \& 300 K undergoes a process by which pressure is doubled. The nature of the process is unspecified but for the process : $\Delta$U = 900 cal. Calculate final volume of gas (in litres) ?\par
[Given: R = 0.08 atm lit./mol/K = 2 cal / K-mol]
\qn{10} Calculate total number of coloured compound (white/colourless not include) out of the following\par
(1) FeS\quad (2) FeS$_2$\quad (3) Zn$_2$[Fe(CN)$_6$]\par
(4) Fe$_3$[Fe(CN)$_6$]$_2$\quad (5) Cu$_2$I$_2$\quad (6) ZnS\quad (7) Ag$_2$S$_2$O$_3$\par
(8) Ag$_2$S\quad (9) ZnO(Hot)\quad (10) CrO$_2$Cl$_2$\quad (11) Cu$_2$[Fe(CN)$_6$]\par
(12) Hg$_2$CrO$_4$\quad (13) KO$_3$\quad (14) KBr\quad (15) H$_2$S
\sect{SECTION-III}
\qn{1} Number of Grignard reagent that can be consumed by given compound.\par
\qfig{p1_chem_s3_q01_1}
\qn{2} Number of compounds which can be oxidised by Tollen's regent.\par
\qfig{p1_chem_s3_q02_1}
\qn{3} The mole fraction of solute in some solution is 1/n. If 50\% of the solute molecules dissociates into two parts and remaining 50\% get dimerized now mole fractions of solvent become 4/5. Find value of n.
\qn{4} An aqueous solution has 2 mole of NH$_4$Cl and this solution is filled between two electrodes which are 20 cm apart. Resistance of this solution is measured to be 100 $\Omega$. Calculate molar conductance (S cm$^2$ mole$^{-1}$) for NH$_4$Cl.
\qn{5} In electrolytic refining of impure copper containing impurities of Ag, Hg, Al, Zn, Au, Fe, Pt, Sn, Mg. Then find out total number of metals present in anode mud in their elemental form.
\qn{6} How many of the following will give N$_2$(g) on heating.\par
(NH$_4$)$_2$SO$_4$, NH$_4$NO$_3$, (NH$_4$)$_2$Cr$_2$O$_7$, NH$_4$Cl, NH$_4$NO$_2$, NH$_4$MnO$_4$.
\qn{7} Find the total stereoisomers for :\par
(a) [Cr(C$_2$O$_4$)$_3$]$^{3-}$ = x\par
(b) [PtCl$_2$(en)$_2$]$^{2+}$ = y\par
(c) [Cr(NH$_3$)$_2$Cl$_2$ (en)]$^+$ = z\par
Hence write the value of (x + y + z).
\qn{8} Anglesite ore + water $\rightarrow$ insoluble\par
Anglesite ore + ammonium acetate solution (aq.) $\xrightarrow{\text{Boiling}}$ Soluble metal containing compound (X) + other products\par
Write the number of oxygen atoms in water soluble metal containing compound X.
\raggedcolumns
\papersec{1}
\partsec{3}{MATHEMATICS}
\sect{SECTION-II (i)}
\qn{1} If f(x) = $\displaystyle\lim_{n\to\infty}\sum_{r=0}^{n}\dfrac{\tan\left(\frac{x}{2^{r+1}}\right)+\tan^3\left(\frac{x}{2^{r+1}}\right)}{\left(1-\tan^2\frac{x}{2^{r+1}}\right)}$, then $\dfrac{\displaystyle\lim_{x\to0}\frac{f(x)}{x}}{2}$ is :-
\qn{2} Let the base of a pyramid is rectangular, three of its vertices of base are A(2, 2, $-$1), B(3, 1, 2) and C(1, 1, 1). Its vertex at the top is P$\left(4,\dfrac{-26}3,\dfrac{-10}3\right)$. Let the coordinates of foot of normal drawn from P on the base of pyramid is ($\ell$, m, n) then $\dfrac{\ell+\mathrm m+\mathrm n}{10}$ is :-
\qn{3} $\dfrac{\displaystyle\sum_{r=1}^{\infty}\tan^{-1}\left(\frac{6.7^{(r-1)}}{1+7^{(2r-1)}}\right)}{\pi}$ is :-
\qn{4} If secA tanB + tanAsecB = 91, then $\dfrac{(\sec\mathrm A\sec\mathrm B+\tan\mathrm A\tan\mathrm B)^2}{4141}$ is
\qn{5} For all complex numbers z$_1$, z$_2$ satisfying $|\mathrm z_1|=12$ and $|\mathrm z_2-(3+4\mathrm i)|=5$, the minimum value of $|\mathrm z_1-\mathrm z_2|$ is :-
\sect{SECTION-II (ii)}
\qn{6} The remainder when 7$^{103}$ is divided by 25 is :-
\qn{7} Let f(x) = max\{x$^2$, (1 $-$ x)$^2$, 2x(1 $-$ x)\}, where 0 $\le$ x $\le$ 1. If the area bounded by y = f(x), x-axis, x = 0 \& x = 1 is $\lambda$, then 27($\lambda$) is :-
\qn{8} There are 10 tickets marked numbers 1,2,3,.....10. Three tickets marked numbers a, b, c (a $>$ b $>$ c) are drawn at random. If the probability that $\displaystyle\int_0^{\mathrm a}\left(\mathrm x^2-2\mathrm{bx}+3\mathrm c\right)\mathrm{dx}=0$ is $\dfrac1{\mathrm p}$, then the value of `p' is
\qn{9} A circle of radius 1 is randomly placed in a 9 $\times$ 12 rectangle ABCD so that the circle lies completely within the rectangle. If $\dfrac{\mathrm m}{\mathrm n}$ (where m and n are relatively prime numbers) denotes the probability that the circle neither touches nor intersect the diagonal AC, then the value of (m + n) is
\qn{10} Let S$_\mathrm n$ be the area of the figure enclosed by a curve y = x$^2$(1 $-$ x)$^\mathrm n$ (0 $\le$ x $\le$ 1) and x-axis, if $\displaystyle\lim_{n\to\infty}\sum_{k=1}^{n}\mathrm S_\mathrm k=\dfrac1{\sqrt{\mathrm A}}$, then `A' is equal to
\sect{SECTION-III}
\qn{1} The sum of $\dfrac3{1.2}\left(\dfrac12\right)+\dfrac4{2.3}\left(\dfrac12\right)^2+\dfrac5{3.4}\left(\dfrac12\right)^3+\ldots\infty$ terms is
\qn{2} Let (1 + x)$^{10}$ = a$_0$ + a$_1$x + a$_2$x$^2$ + ... + a$_{10}$x$^{10}$, then\par
\resizebox{\linewidth}{!}{$\dfrac{(\mathrm a_0-\mathrm a_2+\mathrm a_4-\mathrm a_6+\mathrm a_8-\mathrm a_{10})^2+(\mathrm a_1-\mathrm a_3+\mathrm a_5-\mathrm a_7+\mathrm a_9)^2}{2^{10}}$} is
\qn{3} If f(x + y) = f(x) + f(y) $-$ xy $-$ 1 $\forall$ x, y $\in$ R and f(1) = 1, then number of solutions of f(x) = x, x $\in$ N is
\qn{4} The shortest distance between z-axis and the line x + y + 2z $-$ 3 = 0 = 2x + 3y + 4z $-$ 4 is :-
\qn{5} Let 2ysinx$\dfrac{\mathrm{dy}}{\mathrm{dx}}$ = 2sinxcosx $-$ y$^2$cosx satisfying y$\left(\dfrac\pi2\right)$ = 1, then $\dfrac{\mathrm y^2}{\sin\mathrm x}$ (wherever defined) is :-
\qn{6} $\displaystyle\lim_{x\to0}\dfrac{\left(\frac{\mathrm a^{\mathrm x}+\mathrm b^{\mathrm x}+\mathrm c^{\mathrm x}}{3}\right)^{2/\mathrm x}}{(\mathrm{abc})^{2/3}}$ is :-\par
(Where a, b, c $>$ 0)
\qn{7} Let A = $\begin{bmatrix}0&\alpha\\0&0\end{bmatrix}$ and (A + I)$^{50}$ $-$ 50A = $\begin{bmatrix}\mathrm a&\mathrm b\\\mathrm c&\mathrm d\end{bmatrix}$, then a + b + c + d is
\qn{8} If z$_1$ \& z$_2$ be two distinct complex numbers satisfying $\left|\mathrm z_1^2-\mathrm z_2^2\right|=\left|\bar{\mathrm z}_1^2+\bar{\mathrm z}_2^2-2\bar{\mathrm z}_1\bar{\mathrm z}_2\right|$ \& if (arg z$_1$ $-$ arg z$_2$) = $\dfrac{\mathrm a\pi}{\mathrm b}$, then least possible value of $|$a + b$|$ is, (a,b $\in$ Integer)
\raggedcolumns
\papersec{2}
\partsec{1}{PHYSICS}
\sect{SECTION-I (i)}
\qn{1} Photons of energy 5 eV are incident on cathode. Electrons reaching the anode have kinetic energies varying from 6eV to 8eV.\par
\qfig{p2_phy_s1_q01_1}
\begin{qoptions}
\opt{A}{Work function of the metal is 2 eV.}
\opt{B}{Work function of the metal is 3 eV.}
\opt{C}{Current in the circuit is equal to saturation value.}
\opt{D}{Current in the circuit is less than saturation value.}
\end{qoptions}
\qn{2} We have a massless rod on which two small equally massive balls are attached as in figure (a). The floor and wall are smooth. When the system is released the balls begin to slide as shown in (b). Which of the following quantities are zero when upper ball just loses contact with the wall ?\par
\qfig{p2_phy_s1_q02_1}
\begin{qoptions}
\opt{A}{Tensile force in the rod.}
\opt{B}{Normal between wall \& the ball.}
\opt{C}{Acceleration of the lower ball.}
\opt{D}{Acceleration of the upper ball.}
\end{qoptions}
\qn{3} Two long parallel wires of negligible resistance are connected at one end to a resistance R and at the other end to a ideal d.c. voltage source of emf V. The distance between the axes of the wires is $\eta$ ($\eta\gg1$) times the cross-sectional radius of each wire. The upper wire gets positively charged with linear charge density $\lambda$ and lower wire with linear charge density $-\lambda$. Choose the correct statement(s).\par
\qfig{p2_phy_s1_q03_1}
\begin{qoptions}
\opt{A}{magnitude of linear charge density on wires = $\dfrac{\pi\varepsilon_0\mathrm V}{\ln\eta}$.}
\opt{B}{magnitude of linear charge density on wires = $\dfrac{2\pi\varepsilon_0\mathrm V}{\ln\eta}$.}
\opt{C}{for R = $\sqrt{\dfrac{\mu_0}{\varepsilon_0}\dfrac{\ln\eta}{\pi}}$, the resultant force of interaction between the wires becomes zero.}
\opt{D}{for R = $\sqrt{\dfrac{\mu_0}{\varepsilon_0}}\dfrac{\ln\eta}{\pi}$, the resultant force of interaction between the wires becomes zero.}
\end{qoptions}
\qn{4} AOB is a right angled isosceles triangle with sides OA = OB = $\sqrt2$ a in xy plane. It is carrying current i as shown. The magnetic induction at the point P (0, 0, a)\par
\qfig{p2_phy_s1_q04_1}
\begin{qoptions}
\opt{A}{due to current in AOB is $\vec{\mathrm B}=-\dfrac{\mu_0\mathrm i}{2\sqrt6\,\pi\mathrm a}(\hat{\mathrm i}+\hat{\mathrm j})$}
\opt{B}{due to current in BA is $\vec{\mathrm B}=\dfrac{\mu_0\mathrm i}{4\sqrt6\,\pi\mathrm a}(\hat{\mathrm i}+\hat{\mathrm j}+\sqrt3\hat{\mathrm k})$}
\opt{C}{due to current in AOB is $\vec{\mathrm B}=-\dfrac{\mu_0\mathrm i}{2\sqrt3\,\pi\mathrm a}(\hat{\mathrm i}+\hat{\mathrm j})$}
\opt{D}{due to current in BA is $\vec{\mathrm B}=\dfrac{\mu_0\mathrm i}{4\sqrt6\,\pi\mathrm a}(\hat{\mathrm i}+\hat{\mathrm j}+\sqrt2\hat{\mathrm k})$}
\end{qoptions}
\qn{5} A 5kg block is attached to two springs of force constant 100 N/m each. The block experiences viscous friction. Select the correct statements for its oscillation in steady state. Take $\pi^2$ = 10, b = 4 kg/s.\par
\qfig{p2_phy_s1_q05_1}
\begin{qoptions}
\opt{A}{The velocity amplitude is more than velocity amplitude for any other frequency of applied force.}
\opt{B}{The displacement amplitude is more than displacement amplitude for any other frequency of applied force.}
\opt{C}{Velocity amplitude is 2 m/s}
\opt{D}{Displacement amplitude is $\dfrac2\pi\sqrt{\pi^2-1}$m}
\end{qoptions}
\qn{6} Two identical sleeves A and B each of mass m can slide without friction on a fixed horizontal rod. A load C of mass 2m is suspended from the mid-point of a light inextensible cord of length $\ell$ = 3$\sqrt3$ m connecting the sleeves A and B. Initially the sleeves are at a separation $\ell$ and the load is held motionless. What can you conclude for subsequent motion after the load is released? Acceleration of free fall is 10 m/s$^2$.\par
\qfig{p2_phy_s1_q06_1}
\begin{qoptions}
\opt{A}{Maximum speed of the sleeves is $\sqrt{10}$ m/s.}
\opt{B}{Maximum speed acquired by the load C is 2$\sqrt5$ m/s.}
\opt{C}{The load descends a height 1.5 m to acquire maximum speed.}
\opt{D}{Before the sleeves collide, the net force on the system consisting of the sleeves, the cord and the load becomes zero once.}
\end{qoptions}
\sect{SECTION-I (ii)}
\qn{7} In the circuit shown in figure plates of capacitor C$_2$ are gradually pulled apart. Then match the following :-\par
\qfig{p2_phy_s2_q07_1}
{\footnotesize\setlength{\tabcolsep}{1.5pt}\noindent\begin{tabular}{|c|L{3.7cm}|c|L{2.3cm}|}
\hline
 & \multicolumn{1}{c|}{List-I} & & \multicolumn{1}{c|}{List-II} \\ \hline
(P) & Potential difference across C$_1$ & (1) & Will remain same \\ \hline
(Q) & Energy stored in C$_3$ & (2) & Will increase \\ \hline
(R) & Potential difference across C$_2$ & (3) & Will decrease \\ \hline
(S) & Electric field between the plates of capacitor C$_1$ & (4) & None of these \\ \hline
\end{tabular}}\par
\begin{qoptions}
\opt{A}{P $\rightarrow$ 1;Q $\rightarrow$ 3;R $\rightarrow$ 1;S $\rightarrow$ 4}
\opt{B}{P $\rightarrow$ 3;Q $\rightarrow$ 2;R $\rightarrow$ 4;S $\rightarrow$ 1}
\opt{C}{P $\rightarrow$ 3;Q $\rightarrow$ 1;R $\rightarrow$ 2;S $\rightarrow$ 3}
\opt{D}{P $\rightarrow$ 4;Q $\rightarrow$ 2;R $\rightarrow$ 3;S $\rightarrow$ 4}
\end{qoptions}
\qn{8} List-I shows different charge distributions. List-II gives corresponding electric dipole moments.\par
{\footnotesize\setlength{\tabcolsep}{1.5pt}\noindent\begin{tabular}{|c|L{4.6cm}|c|L{1.8cm}|}
\hline
 & \multicolumn{1}{c|}{List-I} & & \multicolumn{1}{c|}{List-II} \\ \hline
(P) & \qfigi[0.75]{p2_phy_s2_q08_1}\newline Four point charges are kept as shown & (1) & $\sqrt{22}$q$\ell$ \\ \hline
(Q) & \qfigi[0.75]{p2_phy_s2_q08_2}\newline A uniformly charged hollow cone having height $\ell$ and radius $\ell$/3 carries charge $-$q. A charge +q is placed at point `A'. & (2) & $\dfrac{\mathrm q\ell\sqrt2}3$ \\ \hline
(R) & \qfigi[0.75]{p2_phy_s2_q08_3}\newline Uniformly charged arc has charge $-$q. Another charge q is placed at point `A'. Radius of arc is $\ell$ & (3) & $\dfrac{\mathrm q\ell}2$ \\ \hline
(S) & \qfigi[0.75]{p2_phy_s2_q08_4}\newline A cube of side $\ell$ has three charges as shown & (4) & 8$\sqrt2$q$\ell$ \\ \hline
\end{tabular}}\par
\begin{qoptions}
\opt{A}{P $\rightarrow$ 4;Q $\rightarrow$ 2;R $\rightarrow$ 3;S $\rightarrow$ 1}
\opt{B}{P $\rightarrow$ 3;Q $\rightarrow$ 2;R $\rightarrow$ 4;S $\rightarrow$ 4}
\opt{C}{P $\rightarrow$ 3;Q $\rightarrow$ 2;R $\rightarrow$ 4;S $\rightarrow$ 1}
\opt{D}{P $\rightarrow$ 1;Q $\rightarrow$ 4;R $\rightarrow$ 2;S $\rightarrow$ 3}
\end{qoptions}
\qn{9} Match the entries in List-I with all the entries in List-II.\par
{\footnotesize\setlength{\tabcolsep}{1.5pt}\noindent\begin{tabular}{|c|L{5.0cm}|c|c|}
\hline
 & \multicolumn{1}{c|}{List-I} & & List-II \\ \hline
(P) & The current through R$_1$ increase by 0.1 A on adding the lamp as shown below.\newline \qfigi[0.75]{p2_phy_s2_q09_1}\newline If R$_1$/R$_2$ = 4, what is the current through lamp in ampere? & (1) & 0 \\ \hline
(Q) & A uniform thin bar of mass M is sliding along ice while rotating. The velocity of the bar's center of mass is v and its angular velocity is $\omega$. At the moment when the bar is perpendicular to the velocity of center of mass, the bar hits a stationary puck of mass m at a point on the end. The collision is elastic and as a result of collision, bar stays in place and puck slides away.\newline What is the ratio $\dfrac{\mathrm M}{\mathrm m}$ ? & (2) & $\dfrac12$ \\ \hline
(R) & In a meter bridge experiment, when the known resistor R and unknown resistor X are interchanged, the null point shifts by 60 cm. The ratio $\dfrac{\mathrm R}{\mathrm X}$ could be & (3) & 1 \\ \hline
(S) & A long straight wire carrying current i lies parallel to z-axis, along the line y = 0, x = $-$a. If $\vec{\mathrm B}$ be the magnetic field at any point on y-axis, $\displaystyle\int_{-\infty}^{\infty}\vec{\mathrm B}.\overrightarrow{\mathrm{dy}}$ = $\mu_0$ni, find n. & (4) & 4 \\ \hline
\end{tabular}}\par
\begin{qoptions}
\opt{A}{P $\rightarrow$ 4;Q $\rightarrow$ 2;R $\rightarrow$ 3;S $\rightarrow$ 1}
\opt{B}{P $\rightarrow$ 3;Q $\rightarrow$ 2;R $\rightarrow$ 4;S $\rightarrow$ 4}
\opt{C}{P $\rightarrow$ 3;Q $\rightarrow$ 2;R $\rightarrow$ 4;S $\rightarrow$ 1}
\opt{D}{P $\rightarrow$ 2;Q $\rightarrow$ 4;R $\rightarrow$ 4;S $\rightarrow$ 3}
\end{qoptions}
\qn{10} A.C. circuits are shown in list$-$I and i$_{\mathrm{rms}}$ vs $\omega$ graph is shown in list$-$II match the corresponding graph.\par
{\footnotesize\setlength{\tabcolsep}{1.5pt}\noindent\begin{tabular}{|c|L{3.1cm}|c|L{3.1cm}|}
\hline
 & \multicolumn{1}{c|}{List-I} & & \multicolumn{1}{c|}{List-II} \\ \hline
(P) & \qfigi[0.75]{p2_phy_s2_q10_1} & (1) & \qfigi[0.75]{p2_phy_s2_q10_5} \\ \hline
(Q) & \qfigi[0.75]{p2_phy_s2_q10_2} & (2) & \qfigi[0.75]{p2_phy_s2_q10_6} \\ \hline
(R) & \qfigi[0.75]{p2_phy_s2_q10_3} & (3) & \qfigi[0.75]{p2_phy_s2_q10_7} \\ \hline
(S) & \qfigi[0.75]{p2_phy_s2_q10_4} & (4) & \qfigi[0.75]{p2_phy_s2_q10_8} \\ \hline
\end{tabular}}\par
\begin{qoptions}
\opt{A}{P $\rightarrow$ 4;Q $\rightarrow$ 3;R $\rightarrow$ 1;S $\rightarrow$ 2}
\opt{B}{P $\rightarrow$ 3;Q $\rightarrow$ 2;R $\rightarrow$ 1;S $\rightarrow$ 4}
\opt{C}{P $\rightarrow$ 1;Q $\rightarrow$ 2;R $\rightarrow$ 3;S $\rightarrow$ 4}
\opt{D}{P $\rightarrow$ 4;Q $\rightarrow$ 2;R $\rightarrow$ 3;S $\rightarrow$ 1}
\end{qoptions}
\qn{11} A car is taking a turn on a rough horizontal road without slipping as shown in figure. Let F is centripetal force, f the force of friction, N$_1$ and N$_2$ are two normal reactions. As the speed of car is increased match the following two Lists.\par
\qfig{p2_phy_s2_q11_1}
{\footnotesize\setlength{\tabcolsep}{1.5pt}\noindent\begin{tabular}{|c|L{1.6cm}|c|L{3.6cm}|}
\hline
 & \multicolumn{1}{c|}{List-I} & & \multicolumn{1}{c|}{List-II} \\ \hline
(P) & N$_1$ & (1) & Will increase \\ \hline
(Q) & N$_2$ & (2) & Will decrease \\ \hline
(R) & F/f & (3) & Will remains unchanged \\ \hline
(S) & F & (4) & Cannot say anything \\ \hline
\end{tabular}}\par
\begin{qoptions}
\opt{A}{P $\rightarrow$ 1;Q $\rightarrow$ 3;R $\rightarrow$ 2;S $\rightarrow$ 4}
\opt{B}{P $\rightarrow$ 2;Q $\rightarrow$ 1;R $\rightarrow$ 3;S $\rightarrow$ 1}
\opt{C}{P $\rightarrow$ 3;Q $\rightarrow$ 4;R $\rightarrow$ 1;S $\rightarrow$ 4}
\opt{D}{P $\rightarrow$ 4;Q $\rightarrow$ 2;R $\rightarrow$ 4;S $\rightarrow$ 1}
\end{qoptions}
\qn{12} A ball of mass 1kg is projected towards another stationary ball of mass 1kg. Both balls have same radius of 10 cm. If there is a collision, the collision is elastic. List$-$I gives the value of impact parameter and List$-$II gives velocity (magnitude) of ball B after collision. Match them.\par
\qfig{p2_phy_s2_q12_1}
{\footnotesize\setlength{\tabcolsep}{1.5pt}\noindent\begin{tabular}{|c|L{2.2cm}|c|L{2.2cm}|}
\hline
 & \multicolumn{1}{c|}{List-I} & & \multicolumn{1}{c|}{List-II} \\ \hline
(P) & b = 0 & (1) & 10 m/s \\ \hline
(Q) & b = 30 cm & (2) & 0 m/s \\ \hline
(R) & b = 12 cm & (3) & 8 m/s \\ \hline
(S) & b = 16 cm & (4) & 6 m/s \\ \hline
\end{tabular}}\par
\begin{qoptions}
\opt{A}{P $\rightarrow$ 1;Q $\rightarrow$ 2;R $\rightarrow$ 4;S $\rightarrow$ 3}
\opt{B}{P $\rightarrow$ 2;Q $\rightarrow$ 1;R $\rightarrow$ 4;S $\rightarrow$ 3}
\opt{C}{P $\rightarrow$ 1;Q $\rightarrow$ 2;R $\rightarrow$ 3;S $\rightarrow$ 4}
\opt{D}{P $\rightarrow$ 2;Q $\rightarrow$ 1;R $\rightarrow$ 3;S $\rightarrow$ 4}
\end{qoptions}
\sect{SECTION-II}
\qn{1} In the following question x (an unstable nucleus) decays to another nucleus y. In List$-$I different type of decay processes are mentioned. Atomic mass of element of nucleus x and element of nucleus y are M$_\mathrm x$ and M$_\mathrm y$ respectively. The atomic mass of He is M$_{\mathrm{He}}$. The mass of electron is m$_\mathrm e$. Now match the entries of List-I with List-II.\par
{\footnotesize\setlength{\tabcolsep}{1.5pt}\noindent\begin{tabular}{|c|L{3.0cm}|c|L{3.2cm}|}
\hline
 & \multicolumn{1}{c|}{List-I (Decay process)} & & \multicolumn{1}{c|}{List-II (Mass defect)} \\ \hline
(P) & x $\longrightarrow$ y + $\alpha$ & (1) & M$_\mathrm x$ $-$ M$_\mathrm y$ $-$ m$_\mathrm e$ \\ \hline
(Q) & x $\longrightarrow$ y + $_{-1}$e$^0$ & (2) & M$_\mathrm x$ $-$ M$_\mathrm y$ $-$ M$_{\mathrm{He}}$ \\ \hline
(R) & x $\longrightarrow$ y + $_{+1}$e$^0$ & (3) & M$_\mathrm x$ $-$ M$_\mathrm y$ $-$ 2m$_\mathrm e$ \\ \hline
(S) & x + $_{-1}$e$^0$ $\longrightarrow$ y + $\nu$ +X-rays & (4) & M$_\mathrm x$ $-$ M$_\mathrm y$ \\ \hline
\end{tabular}}\par
\qn{2} A white light ray is incident on a glass prism, and it create four refracted rays A, B, C and D. Match the refracted rays with the colours given (1 \& D are rays due to total internal reflection):\par
\qfig{p2_phy_s3_q02_1}
{\footnotesize\setlength{\tabcolsep}{1.5pt}\noindent\begin{tabular}{|c|L{2.2cm}|c|L{2.2cm}|}
\hline
 & \multicolumn{1}{c|}{List-I (Ray)} & & \multicolumn{1}{c|}{List-II (Colour)} \\ \hline
(P) & A & (1) & Red \\ \hline
(Q) & B & (2) & Green \\ \hline
(R) & C & (3) & Yellow \\ \hline
(S) & D & (4) & Blue \\ \hline
\end{tabular}}\par
\qn{3} A polarized electromagnetic plane beam $\vec{\mathrm E}=2\sqrt\pi\left(3\hat{\mathrm i}+4\hat{\mathrm j}\right)$ N/C $\sin\left(\omega\left(\mathrm t+\dfrac{\mathrm Z}{\mathrm C}\right)\right)$ is travelling along $-$ ve z direction. A polaroid is placed in its path at origin so that polaroid is in xy plane. List$-$I indicates pass axis \& List$-$II emergent intensity. Indicated in List$-$I \& List$-$II respectively. Match them\par
{\footnotesize\setlength{\tabcolsep}{1.5pt}\noindent\begin{tabular}{|c|L{2.6cm}|c|L{2.6cm}|}
\hline
 & \multicolumn{1}{c|}{List-I (Pass axis along)} & & \multicolumn{1}{c|}{List-II (Intensity)} \\ \hline
(P) & x = y & (1) & $\dfrac{49}{120}$w/m$^2$ \\ \hline
(Q) & x = 0 & (2) & $\dfrac{48}{125}$w/m$^2$ \\ \hline
(R) & y = 0 & (3) & $\dfrac3{20}$w/m$^2$ \\ \hline
(S) & 3x = 4y & (4) & $\dfrac4{15}$w/m$^2$ \\ \hline
\end{tabular}}\par
\qn{4} A siphon transfers liquid from high level to low level as shown in figure. If $\rho$ is density of liquid, P$_0$ atmospheric pressure, V$_\mathrm a$ is velocity at point a (on the surface of water in the tank) and V$_\mathrm b$ is velocity at point b (at the right open end of the siphon), then (V$_\mathrm b\neq0$). P$_\mathrm a$, P$_\mathrm b$ are pressures at points a and b respectively.\par
\qfig{p2_phy_s3_q04_1}
{\footnotesize\setlength{\tabcolsep}{1.5pt}\noindent\begin{tabular}{|c|L{3.0cm}|c|L{3.0cm}|}
\hline
 & \multicolumn{1}{c|}{List-I} & & \multicolumn{1}{c|}{List-II} \\ \hline
(P) & $(\mathrm h_1+\mathrm h_2)-\dfrac{\mathrm P_0}{\rho\mathrm g}$ & (1) & is equal to one \\ \hline
(Q) & V$_\mathrm a$ $-$ V$_\mathrm b$ & (2) & is less than zero \\ \hline
(R) & h$_1$ & (3) & is greater than zero \\ \hline
(S) & $\dfrac{\mathrm P_\mathrm b-\mathrm P_\mathrm a}{\rho\mathrm g\mathrm h_1}$ & (4) & is equal to zero \\ \hline
\end{tabular}}\par
\qn{5} An ideal gas at initial pressure P$_0$ (= atmospheric pressure) and volume V$_0$ is enclosed by a light piston in a cylinder as shown. Initial the piston is in static equilibrium and the gas is in thermodynamic equilibrium. Now we start to move the piston outwards at a constant slow velocity v. The cross sectional area of piston is A. While the piston is moving we can deliver heat to the gas through a heating filament and pull the piston in such a way that temperature of gas remains constant. Apart from the heat transfer between the gas and the heating filament all other heat exchange can be neglected and the atmospheric pressure is constant. Mark the quantities in List-I with their, appropriate variation in List-II (time t on x-axis)\par
\qfig{p2_phy_s3_q05_1}
{\footnotesize\setlength{\tabcolsep}{1.5pt}\noindent\begin{tabular}{|c|L{3.5cm}|c|L{1.6cm}|}
\hline
 & \multicolumn{1}{c|}{List-I} & & \multicolumn{1}{c|}{List-II} \\ \hline
(P) & Pressure P of gas on y-axis & (1) & \qfigi[0.75]{p2_phy_s3_q05_2} \\ \hline
(Q) & Volume V of gas on y-axis & (2) & \qfigi[0.75]{p2_phy_s3_q05_3} \\ \hline
(R) & Force applied by pulling agent (apart from atmosphere) on the piston on y-axis & (3) & \qfigi[0.75]{p2_phy_s3_q05_4} \\ \hline
(S) & Amount of cumulative work done by gas on piston from start till time `t' on y-axis & (4) & \qfigi[0.75]{p2_phy_s3_q05_5} \\ \hline
\end{tabular}}\par
\qn{6} A plane mirror is tied to the free end of an ideal spring. The other end of the spring is attached to a wall. The spring with mirror is held vertically to the floor can slide along it smoothly. When the spring is at its natural length, the mirror is found to move at speed v with respect to ground. An object moving towards mirror with speed 2v with respect to ground. Match the following : (Assume no collision before mirror comes to rest for 1$^{\mathrm{st}}$ time)\par
\qfig{p2_phy_s3_q06_1}
{\footnotesize\setlength{\tabcolsep}{1.5pt}\noindent\begin{tabular}{|c|L{3.8cm}|c|L{1.2cm}|}
\hline
 & \multicolumn{1}{c|}{List-I} & & \multicolumn{1}{c|}{List-II} \\ \hline
(P) & Speed of image with respect to ground when spring at natural length & (1) & v \\ \hline
(Q) & Speed of image with respect to mirror when spring at natural length & (2) & 0 \\ \hline
(R) & Speed of image with respect to object when spring is at natural length & (3) & 2v \\ \hline
(S) & Speed of image with respect to ground frame when spring at maximum compressed state. & (4) & 3v \\ \hline
\end{tabular}}\par
\raggedcolumns
\qfullbegin
\papersec{2}
\partsec{2}{CHEMISTRY}
\sect{SECTION-I (i)}
\qn{1} \qfig{p2_chem_s1_q01_1}
\qfullsplit
Choose correct option(s) ?
\begin{qoptions}
\opt{A}{(Q) is soluble in aq. NaHCO$_3$ as well as in aq. NaOH}
\opt{B}{Formation of (R) is step growth polymerisation}
\opt{C}{(P) will give yellow crystal with NaOI}
\opt{D}{(R) is polymer terylene}
\end{qoptions}
\qfullend
\qn{2} Which of the reaction(s) will produce 1$^\circ$ amine?
\begin{qoptions}
\opt{A}{\qfigi[0.75]{p2_chem_s1_q02_A}}
\opt{B}{\qfigi[0.75]{p2_chem_s1_q02_B}}
\opt{C}{\qfigi[0.75]{p2_chem_s1_q02_C}}
\opt{D}{\qfigi[0.75]{p2_chem_s1_q02_D}}
\end{qoptions}
\qn{3} Which of the following statement(s) are correct for nitrogen family elements ?
\begin{qoptions}
\opt{A}{Oxidation state from +1 to +4 for nitrogen tend to disproportionate in acidic medium.}
\opt{B}{Phosphorous and arsenic can form d$\pi-$d$\pi$ bond with transition metals.}
\opt{C}{The only basic oxide of N$-$family is Bi$_2$O$_5$.}
\opt{D}{The only trihalide BiF$_3$ is ionic in nature.}
\end{qoptions}
\qn{4} Select the correct statement(s).
\begin{qoptions}
\opt{A}{Compounds of Lanthanoids resemble with alkali metal in properties.}
\opt{B}{Atomic size of lanthanoids decreases regularly on moving left to right.}
\opt{C}{Lanthanoids show +3 oxidation state.}
\opt{D}{Lanthanoids are soft and silvery white.}
\end{qoptions}
\qn{5} Identify correct statement(s) :-
\begin{qoptions}
\opt{A}{Of the three types of cubic lattices, SC, BCC \& FCC those having greatest \& least densities for identical atoms are FCC \& SC respectively.}
\opt{B}{A plot between vapour pressure of a liquid \& its absolute temperature results in a straight line.}
\opt{C}{In CCP three dimensional structure, possible distances between octahedral voids are $\dfrac a2,\dfrac a{\sqrt2}$. [a = Edge length of unit cell]}
\opt{D}{In FCC unit cell the possible distances between tetrahedral voids are $\dfrac a2,\dfrac a{\sqrt2},\dfrac{a\sqrt3}2$ (where a is edge length of unit cell).}
\end{qoptions}
\qn{6} According to the Arrhenius equation:
\begin{qoptions}
\opt{A}{In Arrhenius equation, factor e$^{-\mathrm{Ea/RT}}$ corresponds to fraction of reactant molecules that have kinetic energy equal to activation energy.}
\opt{B}{Rate constant increases with increases in temperature. This is mainly due to a greater number of collisions whose energy exceeds the activation energy.}
\opt{C}{Higher the magnitude of activation energy of a reaction, stronger is the temperature dependence of its rate constant.}
\opt{D}{The pre$-$exponential factor depends on the frequency at which collisions between reactant molecules per unit volume per unit time occur, irrespective of their energy.}
\end{qoptions}
\sect{SECTION-I (ii)}
\qn{7} Matching List.\par
{\footnotesize\setlength{\tabcolsep}{1.5pt}\noindent\begin{tabular}{|c|L{3.9cm}|c|L{2.4cm}|}
\hline
 & \multicolumn{1}{>{\centering\arraybackslash}p{3.9cm}|}{List-I} & & \multicolumn{1}{>{\centering\arraybackslash}p{2.4cm}|}{List-II (About product mixture)} \\ \hline
(P) & \qfigi[0.75]{p2_chem_s2_q07_1} & (1) & Optically active \\ \hline
(Q) & \qfigi[0.75]{p2_chem_s2_q07_2} & (2) & Diastereoisomers \\ \hline
(R) & \qfigi[0.75]{p2_chem_s2_q07_3} & (3) & Optically inactive \\ \hline
(S) & \qfigi[0.75]{p2_chem_s2_q07_4} & (4) & Meso \\ \hline
\end{tabular}}\par
\begin{qoptions}
\opt{A}{P $\rightarrow$ 1,2;Q $\rightarrow$ 3;R $\rightarrow$ 3;S $\rightarrow$ 3,4}
\opt{B}{P $\rightarrow$ 1,3;Q $\rightarrow$ 2;R $\rightarrow$ 3;S $\rightarrow$ 3,4}
\opt{C}{P $\rightarrow$ 1,2;Q $\rightarrow$ 3;R $\rightarrow$ 1;S $\rightarrow$ 2,3}
\opt{D}{P $\rightarrow$ 1,2;Q $\rightarrow$ 4;R $\rightarrow$ 2;S $\rightarrow$ 3,4}
\end{qoptions}
\qn{8} Match the list\par
{\footnotesize\setlength{\tabcolsep}{1.5pt}\noindent\begin{tabular}{|c|L{4.6cm}|c|L{2.0cm}|}
\hline
 & \multicolumn{1}{>{\centering\arraybackslash}p{4.6cm}|}{List-I (Reactions)} & & \multicolumn{1}{>{\centering\arraybackslash}p{2.0cm}|}{List-II (A and B is differentiated by)} \\ \hline
(P) & \qfigi[0.75]{p2_chem_s2_q08_1} & (1) & FeCl$_3$ \\ \hline
(Q) & \qfigi[0.75]{p2_chem_s2_q08_2}\newline \qfigi[0.75]{p2_chem_s2_q08_3} & (2) & Fehling test \\ \hline
(R) & \qfigi[0.75]{p2_chem_s2_q08_4} & (3) & NaHCO$_3$ \\ \hline
(S) & \qfigi[0.75]{p2_chem_s2_q08_5} & (4) & Iodoform test \\ \hline
\end{tabular}}\par
\begin{qoptions}
\opt{A}{P $\rightarrow$ 3;Q $\rightarrow$ 2;R $\rightarrow$ 4;S $\rightarrow$ 2}
\opt{B}{P $\rightarrow$ 2;Q $\rightarrow$ 3;R $\rightarrow$ 1;S $\rightarrow$ 3}
\opt{C}{P $\rightarrow$ 2;Q $\rightarrow$ 3;R $\rightarrow$ 3;S $\rightarrow$ 1}
\opt{D}{P $\rightarrow$ 3;Q $\rightarrow$ 2;R $\rightarrow$ 4;S $\rightarrow$ 4}
\end{qoptions}
\qn{9} Match list-I with list-II.\par
{\footnotesize\setlength{\tabcolsep}{1.5pt}\noindent\begin{tabular}{|c|L{4.6cm}|c|L{1.2cm}|}
\hline
 & \multicolumn{1}{>{\centering\arraybackslash}p{4.6cm}|}{List-I [Property of a metal (M)]} & & \multicolumn{1}{>{\centering\arraybackslash}p{1.2cm}|}{List-II [Metal (M)]} \\ \hline
(P) & "M" produces H$_2$ with KOH solution & (1) & Au \\ \hline
(Q) & "M" produces M$^{2+}$ ion when reacts with dil. HCl & (2) & Zn \\ \hline
(R) & "M" does not dissolve in dil. HCl but dissolves in dil. HNO$_3$ & (3) & Hg \\ \hline
(S) & "M" does not dissolve in conc. HNO$_3$ dil. HNO$_3$ or H$_2$SO$_4$ & (4) & Fe \\ \hline
\end{tabular}}\par
\begin{qoptions}
\opt{A}{P $\rightarrow$ 1;Q $\rightarrow$ 2;R $\rightarrow$ 3;S $\rightarrow$ 2}
\opt{B}{P $\rightarrow$ 4;Q $\rightarrow$ 1;R $\rightarrow$ 2;S $\rightarrow$ 3}
\opt{C}{P $\rightarrow$ 2;Q $\rightarrow$ 4;R $\rightarrow$ 1;S $\rightarrow$ 4}
\opt{D}{P $\rightarrow$ 2;Q $\rightarrow$ 2;R $\rightarrow$ 4;S $\rightarrow$ 3}
\end{qoptions}
\qn{10} Match list-I with list-II.\par
{\footnotesize\setlength{\tabcolsep}{1.5pt}\noindent\begin{tabular}{|c|L{4.6cm}|c|L{2.6cm}|}
\hline
 & \multicolumn{1}{>{\centering\arraybackslash}p{4.6cm}|}{List-I} & & \multicolumn{1}{>{\centering\arraybackslash}p{2.6cm}|}{List-II} \\ \hline
(P) & K$_4$[Fe(CN)$_6$] $>$ K$_3$[Fe(CN)$_6$] & (1) & Correct order of stability \\ \hline
(Q) & [Co(NH$_3$)$_6$]Cl$_3$ $<$ [Co(en)$_3$]Cl$_3$ & (2) & Correct order of magnetic moment \\ \hline
(R) & [Co(NH$_3$)$_6$]Cl$_3$ $<$ Na$_3$[CoF$_6$] & (3) & Correct order of oxidation number of central metal atom \\ \hline
(S) & [NiF$_6$]$^{2-}$ $>$ [NiF$_6$]$^{4-}$ & (4) & Correct order of number of t$_{2g}$ electrons \\ \hline
\end{tabular}}\par
\begin{qoptions}
\opt{A}{P $\rightarrow$ 4;Q $\rightarrow$ 1;R $\rightarrow$ 2;S $\rightarrow$ 1,3}
\opt{B}{P $\rightarrow$ 3;Q $\rightarrow$ 3,4;R $\rightarrow$ 1,3,4;S $\rightarrow$ 3,4}
\opt{C}{P $\rightarrow$ 2,3;Q $\rightarrow$ 2,4;R $\rightarrow$ 4;S $\rightarrow$ 3,4}
\opt{D}{P $\rightarrow$ 3,4;Q $\rightarrow$ 2,4;R $\rightarrow$ 4;S $\rightarrow$ 2,3,4}
\end{qoptions}
\qn{11} \par
{\footnotesize\setlength{\tabcolsep}{1.5pt}\noindent\begin{tabular}{|c|L{4.2cm}|c|L{2.6cm}|}
\hline
 & \multicolumn{1}{>{\centering\arraybackslash}p{4.2cm}|}{List-I} & & \multicolumn{1}{>{\centering\arraybackslash}p{2.6cm}|}{List-II} \\ \hline
(P) & Orbital angular momentum of an electron & (1) & $\sqrt{s(s+1)}\,\dfrac{h}{2\pi}$ \\ \hline
(Q) & Angular momentum of an electron in an orbit & (2) & $\sqrt{n(n+2)}$ \\ \hline
(R) & Spin angular momentum of an electron & (3) & $\dfrac{nh}{2\pi}$ \\ \hline
(S) & Spin magnetic moment of atom/ion & (4) & $\sqrt{\ell(\ell+1)}\,\dfrac{h}{2\pi}$ \\ \hline
\end{tabular}}\par
\begin{qoptions}
\opt{A}{P $\rightarrow$ 1;Q $\rightarrow$ 2;R $\rightarrow$ 3;S $\rightarrow$ 4}
\opt{B}{P $\rightarrow$ 2;Q $\rightarrow$ 3;R $\rightarrow$ 1;S $\rightarrow$ 4}
\opt{C}{P $\rightarrow$ 4;Q $\rightarrow$ 3;R $\rightarrow$ 1;S $\rightarrow$ 2}
\opt{D}{P $\rightarrow$ 3;Q $\rightarrow$ 4;R $\rightarrow$ 2;S $\rightarrow$ 1}
\end{qoptions}
\qn{12} Match list-I with list-II. (Assuming all the solutes are non volatile and all solutions are ideal.)\par
{\footnotesize\setlength{\tabcolsep}{1.5pt}\noindent\begin{tabular}{|c|L{4.6cm}|c|L{2.4cm}|}
\hline
 & \multicolumn{1}{>{\centering\arraybackslash}p{4.6cm}|}{List - I (Experiment)} & & \multicolumn{1}{>{\centering\arraybackslash}p{2.4cm}|}{List - II (Result)} \\ \hline
(P) & 10 milli moles of HgI$_2$ are added to 10 ml of 2 M aqueous NaI & (1) & Osmotic pressure of solution increases \\ \hline
(Q) & 10 milli moles of KCl are added to 10 ml of 1 M aqueous NaCl & (2) & Vapour pressure of solution increases \\ \hline
(R) & 10 ml 0.1 M NaOH aqueous solution is added to 10 ml 0.1 M HCl aqueous solution & (3) & Boiling point of solution increases \\ \hline
(S) & 10 ml of 0.1 M NaOH aqueous solution is added to 10 ml 0.1 M CH$_3$COOH aqueous solution & (4) & Freezing point of solution increases \\ \hline
\end{tabular}}\par
\begin{qoptions}
\opt{A}{P $\rightarrow$ 1;Q $\rightarrow$ 4;R $\rightarrow$ 2;S $\rightarrow$ 4}
\opt{B}{P $\rightarrow$ 2;Q $\rightarrow$ 4;R $\rightarrow$ 1;S $\rightarrow$ 3}
\opt{C}{P $\rightarrow$ 1;Q $\rightarrow$ 2;R $\rightarrow$ 1;S $\rightarrow$ 2}
\opt{D}{P $\rightarrow$ 2;Q $\rightarrow$ 4;R $\rightarrow$ 2;S $\rightarrow$ 4}
\end{qoptions}
\sect{SECTION-II}
\qn{1} \par
{\footnotesize\setlength{\tabcolsep}{1.5pt}\noindent\begin{tabular}{|c|L{4.2cm}|c|L{2.6cm}|}
\hline
 & \multicolumn{1}{>{\centering\arraybackslash}p{4.2cm}|}{List-I (Compound)} & & \multicolumn{1}{>{\centering\arraybackslash}p{2.6cm}|}{List-II (Characteristic)} \\ \hline
(P) & \qfigi[0.75]{p2_chem_s3_pre_1} & (1) & A derivative of dipeptide \\ \hline
(Q) & \qfigi[0.75]{p2_chem_s3_pre_2} & (2) & A derivative of disaccharide \\ \hline
(R) & \qfigi[0.75]{p2_chem_s3_pre_3} & (3) & Reduction product of an aldopentose \\ \hline
(S) & \qfigi[0.75]{p2_chem_s3_pre_4} & (4) & Oxidation product of aldohexose \\ \hline
\end{tabular}}\par
Match each process given in List-I with one or more effect(s) in List-II.\par
\qn{2} \par
{\footnotesize\setlength{\tabcolsep}{1.5pt}\noindent\begin{tabular}{|c|L{5.4cm}|c|L{1.9cm}|}
\hline
 & \multicolumn{1}{>{\centering\arraybackslash}p{5.4cm}|}{List-I (Reactant)} & & \multicolumn{1}{>{\centering\arraybackslash}p{1.9cm}|}{List-II (Product)} \\ \hline
(P) & \qfigi[0.75]{p2_chem_s3_q01_1} & (1) & 1,3-Dimethyl benzene \\ \hline
(Q) & \qfigi[0.75]{p2_chem_s3_q01_4} & (2) & Toluene \\ \hline
(R) & \qfigi[0.75]{p2_chem_s3_q01_2} & (3) & 1,2-dimethyl benzene \\ \hline
(S) & \qfigi[0.75]{p2_chem_s3_q01_3} & (4) & 1,3,5-Trimethyl benzene \\ \hline
\end{tabular}}\par
Match each process given in List-I with one or more effect(s) in List-II.\par
\qn{3} \par
{\footnotesize\setlength{\tabcolsep}{1.5pt}\noindent\begin{tabular}{|c|L{3.4cm}|c|L{2.0cm}|}
\hline
 & \multicolumn{1}{>{\centering\arraybackslash}p{3.4cm}|}{List-I (Name of mineral)} & & \multicolumn{1}{>{\centering\arraybackslash}p{2.0cm}|}{List-II (Metal present in given mineral)} \\ \hline
(P) & Fluorapatite & (1) & Zn \\ \hline
(Q) & Kaolinite & (2) & Al \\ \hline
(R) & Magnesite & (3) & Ca \\ \hline
(S) & Calamine & (4) & Mg \\ \hline
\end{tabular}}\par
\qn{4} \par
{\footnotesize\setlength{\tabcolsep}{1.5pt}\noindent\begin{tabular}{|c|L{3.9cm}|c|L{2.4cm}|}
\hline
 & \multicolumn{1}{>{\centering\arraybackslash}p{3.9cm}|}{List-I (Reaction)} & & \multicolumn{1}{>{\centering\arraybackslash}p{2.4cm}|}{List-II (colour of metal containing product)} \\ \hline
(P) & MnO$_2$ + KOH $\xrightarrow{\mathrm{O_2}}$ & (1) & Deep blue \\ \hline
(Q) & Co$^{2+}$ + KNO$_2$ $\xrightarrow{\mathrm{H^+}}$ & (2) & Yellow \\ \hline
(R) & Mn$^{+2}$ + S$_2$O$_8^{2-}$ + H$_2$O $\rightarrow$ & (3) & Dark green \\ \hline
(S) & Ni(OH)$_2$ + NH$_4^+$ + OH$^-$ $\rightarrow$ & (4) & Purple \\ \hline
\end{tabular}}\par
\qn{5} \par
{\footnotesize\setlength{\tabcolsep}{1.5pt}\noindent\begin{tabular}{|c|L{3.9cm}|c|L{2.4cm}|}
\hline
 & \multicolumn{1}{>{\centering\arraybackslash}p{3.9cm}|}{List-I (For a sample of ideal gas)} & & \multicolumn{1}{>{\centering\arraybackslash}p{2.4cm}|}{List-II} \\ \hline
(P) & $\Delta$S$_{\mathrm{Sys}}$ for isobaric process & (1) & nC$_\mathrm{V}\,\ell$n$\dfrac{\mathrm T_2}{\mathrm T_1}$ \\ \hline
(Q) & $\Delta$S$_{\mathrm{Sys}}$ for isothermal process & (2) & nRT $\ell$n$\dfrac{\mathrm V_1}{\mathrm V_2}$ \\ \hline
(R) & Work done in reversible isothermal process & (3) & $-$nR $\ell$n$\dfrac{\mathrm V_1}{\mathrm V_2}$ \\ \hline
(S) & $\Delta$S$_{\mathrm{Sys}}$ for isochoric process & (4) & nC$_\mathrm{P}\,\ell$n$\dfrac{\mathrm T_2}{\mathrm T_1}$ \\ \hline
\end{tabular}}\par
\qn{6} Match the Lists.\par
{\footnotesize\setlength{\tabcolsep}{1.5pt}\noindent\begin{tabular}{|c|L{4.9cm}|c|L{1.5cm}|}
\hline
 & \multicolumn{1}{>{\centering\arraybackslash}p{4.9cm}|}{List-I} & & \multicolumn{1}{>{\centering\arraybackslash}p{1.5cm}|}{List-II (pH)} \\ \hline
(P) & 2.5 ml $\dfrac25$ M weak base (K$_\mathrm b$ = 10$^{-10}$) is mixed with $\dfrac2{15}$ M HCl for complete neutralisation & (1) & 5 \\ \hline
(Q) & 1 mole Na$_3$A is mixed with 2.5 mole of HCl to give 2 L solution (pKa$_1$, pKa$_2$, pKa$_3$ of H$_3$A are 3, 6, 10 respectively) & (2) & 2.5 \\ \hline
(R) & pH at which 0.1 M CH$_3$COOH dissociates $\dfrac{1000}{11}$\% (pKa of CH$_3$COOH = 5) & (3) & 3 \\ \hline
(S) & 1 L 0.1 M HA is mixed with 99 L water (pKa of HA = pKb of A$^-$) & (4) & 6 \\ \hline
\end{tabular}}\par
\raggedcolumns
\papersec{2}
\partsec{3}{MATHEMATICS}
\sect{SECTION-I (i)}
\qn{1} Let\par
\resizebox{\linewidth}{!}{$\mathrm f(x)=\displaystyle\lim_{n\to\infty}(-n^3)\left(\left|2\tan^{-1}\left(\mathrm e^{|x|}-\mathrm e\right)-\dfrac1{n^3}\right|-2\left|\tan^{-1}\left(\mathrm e^{|x|}-\mathrm e\right)\right|\right),\ x\in\mathbb R$}\par
then which of the following statement(s) is(are) correct ?
\begin{qoptions}
\opt{A}{The number of points where y = f(x) is discontinuous is 1}
\opt{B}{The number of points where y = $|$f(x)$|$ is discontinuous is 1}
\opt{C}{f$\left(\dfrac12\right)$ + f(3) = 0}
\opt{D}{Number of solution of f(x) = sgn(x) is infinite. (Where sgn(x) is the signum function of x)}
\end{qoptions}
\qn{2} If f : R $\rightarrow$ R be differentiable function such that (f(x))$^7$ = x $-$ f(x), then
\begin{qoptions}
\opt{A}{$\displaystyle\int_0^{\sqrt2}\left(\mathrm f^{-1}(\mathrm x)\right)\mathrm{dx}=3$}
\opt{B}{f(x) is increasing $\forall$ x $\in$ R}
\opt{C}{$\dfrac{\mathrm d}{\mathrm{dx}}\left(\mathrm f^{-1}(\mathrm x)\right)$ at x = 2 is 499}
\opt{D}{Let g(x) be the inverse of f(x). Then g$''$(1) = 42}
\end{qoptions}
\qn{3} P denotes plane containing all points equidistant from A($-$4, 2, 1), B(2, $-$4, 3). Q be the plane x $-$ y + cz = 1, where c $\in$ R, then
\begin{qoptions}
\opt{A}{If the line L with equation $\dfrac{\mathrm x-1}1=\dfrac{\mathrm y+2}3=\dfrac{\mathrm z-7}{-1}$ intersect the plane P at (x$_0$, y$_0$, z$_0$) then x$_0$ + y$_0$ + z$_0$ = 12}
\opt{B}{If angle between the planes P and Q is 45$^\circ$ then product of all values of c is $-$2}
\opt{C}{The plane P is parallel to plane Q iff c = $\dfrac13$}
\opt{D}{The plane P is parallel to plane Q iff c = 1}
\end{qoptions}
\qn{4} In a drawer, Mr. Abhishek has 5 pairs of socks, each pair having a different colour. On Monday, Mr. Abhishek selects two socks at random from the 10 socks in the drawer. On Tuesday, Mr. Abhishek selects 2 of the remaining 8 socks at random; and on Wednesday, two of the remaining 6 socks at random. The probability that Wednesday is the first day Mr. Abhishek select same colour socks is $\dfrac mn$, where m and n are relatively prime positive integers, then
\begin{qoptions}
\opt{A}{m + n = 341}
\opt{B}{$\displaystyle\lim_{x\to0^+}x\left(\left[\dfrac1x\right]+\left[\dfrac2x\right]+\left[\dfrac3x\right]+....+\left[\dfrac mx\right]\right)$ = m + n ([.] is greates integer function)}
\opt{C}{Number of ordered pairs (x,y), (x and y are natural numbers), satisfying the condition x + y $<$ 28 and x + y $\ge$ 6, is m + n}
\opt{D}{$\displaystyle\lim_{x\to0}x\left(\left[\dfrac1x\right]+\left[\dfrac2x\right]+\left[\dfrac3x\right]+...+\left[\dfrac mx\right]\right)$ = n $-$ m ([.] is greates integer function)}
\end{qoptions}
\qn{5} A line divides a plane into 2 regions. Two lines divide the plane into maximum 4 regions. If L$_n$ is the maximum number of regions divided by n lines then which of the following is/are true :
\begin{qoptions}
\opt{A}{L$_{20}$ = 211}
\opt{B}{L$_{10}$ = 56}
\opt{C}{L$_{15}$ = 121}
\opt{D}{L$_{25}$ = 326}
\end{qoptions}
\qn{6} If A $=\begin{pmatrix}2&0&1\\3&2&0\\1&0&1\end{pmatrix}$, then which of the following is(are) INCORRECT (where A$^{-1}$ is inverse of matrix A, I is an unit matrix of order 3, adj A is adjoint of matrix A, $|$A$|$ denotes determinant value of matrix A.)
\begin{qoptions}
\opt{A}{$\dfrac1{81}\left|\mathrm{adj}(2\mathrm A)\cdot\mathrm{adj}\left(3\mathrm A^{-1}\right)\cdot\mathrm A^{-1}\right|$ is equal to 288}
\opt{B}{$\left(\mathrm A^{-1}\right)^2=\dfrac12\left(\mathrm A-5\mathrm I+7\mathrm A^{-1}\right)$}
\opt{C}{$\left|\mathrm A^5-2\mathrm A^4\right|$ is equal to 16}
\opt{D}{A$^{-1}=\dfrac12\left(\mathrm A^2-5\mathrm A+6\mathrm I\right)$}
\end{qoptions}
\sect{SECTION-I (ii)}
\qn{7} Match the following :\par
{\footnotesize\setlength{\tabcolsep}{1.5pt}\noindent\begin{tabular}{|c|L{5.6cm}|c|L{0.6cm}|}
\hline
 & \multicolumn{1}{>{\centering\arraybackslash}p{5.6cm}|}{List-I} & & \multicolumn{1}{>{\centering\arraybackslash}p{0.6cm}|}{List-II} \\ \hline
(P) & If the distance of the plane passing through the point P(1,1,1) and perpendicular to the line $\dfrac{x-1}3=\dfrac{y-1}0=\dfrac{z-1}4$ from the origin is $\dfrac pq$ (where p and q are coprime), then the value of (p $-$ q) is & (1) & 1 \\ \hline
(Q) & The interval [0,4] is divided into n equal sub-intervals by the points x$_0$ = 0, x$_1$,x$_2$.......x$_{n-1}$, x$_n$ = 4. If $\delta$x = x$_i$ $-$ x$_{i-1}$ for i = 1,2,3,... then the value of $\displaystyle\lim_{\delta x\to0}\sum_{i=1}^{n-1}x_i\,\delta x$ is & (2) & 2 \\ \hline
(R) & Let f : R $\rightarrow$ R be a function satisfying xf(x) + (1 $-$ x) f($-$x) = x$^2$ + x + 1 for any real numbers x. The greatest real number M for which f(x) $\ge$ M for all real numbers x, is equal to $\dfrac pq$, where p and q are coprime. The value of (q $-$ p) is & (3) & 3 \\ \hline
(S) & The area of the trapezium ABCD with AB$\|$CD, AD $\perp$ AB and AB = 3CD is equal to 4. A circle inside the trapezium is tangent to all of its sides. If the radius of the circle is r then the value of 4r$^2$, is & (4) & 8 \\ \hline
\end{tabular}}\par
\begin{qoptions}
\opt{A}{P $\rightarrow$ 4;Q $\rightarrow$ 2;R $\rightarrow$ 4;S $\rightarrow$ 1}
\opt{B}{P $\rightarrow$ 3;Q $\rightarrow$ 1;R $\rightarrow$ 2;S $\rightarrow$ 4}
\opt{C}{P $\rightarrow$ 2;Q $\rightarrow$ 4;R $\rightarrow$ 1;S $\rightarrow$ 3}
\opt{D}{P $\rightarrow$ 1;Q $\rightarrow$ 2;R $\rightarrow$ 3;S $\rightarrow$ 4}
\end{qoptions}
\qn{8} Match the following :\par
{\footnotesize\setlength{\tabcolsep}{1.5pt}\noindent\begin{tabular}{|c|L{5.6cm}|c|L{0.6cm}|}
\hline
 & \multicolumn{1}{>{\centering\arraybackslash}p{5.6cm}|}{List-I} & & \multicolumn{1}{>{\centering\arraybackslash}p{0.6cm}|}{List-II} \\ \hline
(P) & Total number of values of x, of the form $\dfrac1n$, n $\in$ N in the interval $\left[\dfrac1{15},\dfrac1{10}\right]$ which satisfies \{x\} + \{2x\} +...+ \{12x\} = 78x where \{ \} denotes fractional part of x & (1) & 1 \\ \hline
(Q) & Let (2x$^2$ + 3x + 4)$^{10}$ = $\displaystyle\sum_{r=0}^{20}\mathrm a_r\mathrm x^r$ then the value of $\dfrac{2\mathrm a_9}{\mathrm a_{11}}$ is equal to & (2) & 2 \\ \hline
(R) & Let G be centroid of $\Delta$ABC whose side length are a,b,c. If P is a point in the plane of $\Delta$ABC such that PA = 1, PB = 2, PC = 1 \& PG = 1 then the value of $\dfrac{a^2+b^2+c^2}9$ is equal to & (3) & 3 \\ \hline
(S) & If a derivable function f : R$^+$ $\rightarrow$ R satisfies f(xy) = f(x) + f(y) $\forall$ x,y $\in$ R$^+$. If f(16) = 3 then $\dfrac1{\mathrm f(2)}+\dfrac1{\mathrm f(4)}$ is equal to & (4) & 4 \\ \hline
\end{tabular}}\par
\begin{qoptions}
\opt{A}{P $\rightarrow$ 3;Q $\rightarrow$ 4;R $\rightarrow$ 1;S $\rightarrow$ 2}
\opt{B}{P $\rightarrow$ 2;Q $\rightarrow$ 4;R $\rightarrow$ 1;S $\rightarrow$ 3}
\opt{C}{P $\rightarrow$ 4;Q $\rightarrow$ 3;R $\rightarrow$ 1;S $\rightarrow$ 2}
\opt{D}{P $\rightarrow$ 4;Q $\rightarrow$ 1;R $\rightarrow$ 2;S $\rightarrow$ 3}
\end{qoptions}
\qn{9} Match the following :\par
{\footnotesize\setlength{\tabcolsep}{1.5pt}\noindent\begin{tabular}{|c|L{5.6cm}|c|L{0.6cm}|}
\hline
 & \multicolumn{1}{>{\centering\arraybackslash}p{5.6cm}|}{List-I} & & \multicolumn{1}{>{\centering\arraybackslash}p{0.6cm}|}{List-II} \\ \hline
(P) & If the distance between two parallel tangents having slope m drawn to the hyperbola $\dfrac{x^2}9-\dfrac{y^2}{49}=1$ is 2, then $\dfrac25|$m$|$ is equal to & (1) & 5 \\ \hline
(Q) & If a variable line has its intercepts on the coordiante axes e and e$'$, where $\dfrac e2$ and $\dfrac{e'}2$ are the eccentricities of a hyperbola and its conjugate hyperbola, then the lines always touches the circle x$^2$ + y$^2$ = r$^2$, where r is equal to & (2) & 4 \\ \hline
(R) & If the equation of lines touching both parabolas y$^2$ $-$ 4x = 0 and x$^2$ + 32y = 0 is ax $-$ 2y + b = 0 (a,b $\in$ R), then (a + b) is equal to & (3) & 2 \\ \hline
(S) & If the mid-point of a chord of the ellipse $\dfrac{x^2}{16}+\dfrac{y^2}{25}=1$ is M(0,3) and the length of the chord is $\dfrac{8p}5$, then p is equal to & (4) & 1 \\ \hline
\end{tabular}}\par
\begin{qoptions}
\opt{A}{P $\rightarrow$ 3;Q $\rightarrow$ 2;R $\rightarrow$ 1;S $\rightarrow$ 4}
\opt{B}{P $\rightarrow$ 3;Q $\rightarrow$ 4;R $\rightarrow$ 2;S $\rightarrow$ 1}
\opt{C}{P $\rightarrow$ 2;Q $\rightarrow$ 3;R $\rightarrow$ 4;S $\rightarrow$ 1}
\opt{D}{P $\rightarrow$ 4;Q $\rightarrow$ 3;R $\rightarrow$ 1;S $\rightarrow$ 2}
\end{qoptions}
\qn{10} Match the following List-I with List-II\par
{\footnotesize\setlength{\tabcolsep}{1.5pt}\noindent\begin{tabular}{|c|L{6.0cm}|c|L{0.6cm}|}
\hline
 & \multicolumn{1}{>{\centering\arraybackslash}p{6.0cm}|}{List-I} & & \multicolumn{1}{>{\centering\arraybackslash}p{0.6cm}|}{List-II} \\ \hline
(P) & Let f(x) = x$^2$ sin$\dfrac\pi x$ (x $\ne$ 0), then the number of points in $\left(\dfrac1{21},1\right)$ at which f$'$(x) = 0 is & (1) & 18 \\ \hline
(Q) & Let f(x) = x$^2$ + px + q + 3, if f(f(1)) = f(f($-$2)) = 0 where f(1) $\ne$ f($-$2) and g(x) = x$^2$ + (8p$^2$ $-$ 18p)x $-$ q + 2, then 2g(2) is & (2) & 20 \\ \hline
(R) & Let f(x) = x$^9$ $-$ 3x$^8$ $-$ 53x$^7$ $-$ 8x$^6$ $-$ 8x$^5$ $-$ 8x$^4$ $-$ 8x$^3$ $-$ 8x$^2$ $-$ 8x + 9, then f(9) will be equal to & (3) & 23 \\ \hline
(S) & Let L = \resizebox{5.8cm}{!}{$\displaystyle\lim_{x\to0}\left(1+\mathrm e^{-\frac1{x^2}}\tan^{-1}\dfrac1{x^2}+x^2\mathrm e^{-\frac1{x^2}}\sin\dfrac1{x^4}\right)^{\mathrm e^{\left(\frac1{x^2}\right)}}$}, then [15 log$_\mathrm e$L] is (where [.] represents greater integer function) & (4) & 27 \\ \hline
\end{tabular}}\par
The correct option is :
\begin{qoptions}
\opt{A}{P $\rightarrow$ 2;Q $\rightarrow$ 4;R $\rightarrow$ 1;S $\rightarrow$ 3}
\opt{B}{P $\rightarrow$ 4;Q $\rightarrow$ 3;R $\rightarrow$ 2;S $\rightarrow$ 1}
\opt{C}{P $\rightarrow$ 4;Q $\rightarrow$ 3;R $\rightarrow$ 1;S $\rightarrow$ 2}
\opt{D}{P $\rightarrow$ 2;Q $\rightarrow$ 4;R $\rightarrow$ 2;S $\rightarrow$ 3}
\end{qoptions}
\qn{11} Match the following List-I with List-II\par
{\footnotesize\setlength{\tabcolsep}{1.5pt}\noindent\begin{tabular}{|c|L{6.0cm}|c|L{0.6cm}|}
\hline
 & \multicolumn{1}{>{\centering\arraybackslash}p{6.0cm}|}{List-I} & & \multicolumn{1}{>{\centering\arraybackslash}p{0.6cm}|}{List-II} \\ \hline
(P) & A continuous and differentiable function f(x) satisfies functional equation f(x + y) = $\dfrac{\mathrm f(x)\,\mathrm f(y)}{\mathrm f(x)+\mathrm f(y)}$; f(1) = 16, then f(4) equals & (1) & $\dfrac7{12}$ \\ \hline
(Q) & Let f : R $\rightarrow$ R be a differentiable function satisfying f$\left(\dfrac{x+y}5\right)=\dfrac{6+\mathrm f(x)+\mathrm f(y)}5$ $\forall$ x, y $\in$ R and f$'$(2) = 2, then area bounded by relation $|$y$|$ = h(x), x = 2 and x = $-$2 where h(x) = $|$f$|$x$||$ $-$ 4$|$ is & (2) & 4 \\ \hline
(R) & Let f be a differentiable function satisfying the relation f(x + 2y) = f(x) + f(2y) + e$^{x+2y}$(x + 2y) $-$ xe$^x$ $-$ 2ye$^{2y}$ + 4xy $\forall$ x,y $\in$ R and f$'$(0) = 1, then the value of $\displaystyle\lim_{x\to0}\dfrac{\mathrm f(2x)-4x^2-2x}{x^2}$ is & (3) & 2 \\ \hline
(S) & Let f(x) is a differentiable function satisfying f(y) $-$ f(x) = $\dfrac{x^x}{y^y}$f$\left(\dfrac{y^y}{x^x}\right)$ $\forall$ x,y $\in$ R$^+$, if f$'$(1) = 1, then area bounded between the curve f(x) and g(x) = 2x $-$ 2x$^2$ is & (4) & 8 \\ \hline
\end{tabular}}\par
The correct option is :
\begin{qoptions}
\opt{A}{P $\rightarrow$ 2;Q $\rightarrow$ 3;R $\rightarrow$ 4;S $\rightarrow$ 1}
\opt{B}{P $\rightarrow$ 2;Q $\rightarrow$ 4;R $\rightarrow$ 2;S $\rightarrow$ 1}
\opt{C}{P $\rightarrow$ 4;Q $\rightarrow$ 3;R $\rightarrow$ 1;S $\rightarrow$ 2}
\opt{D}{P $\rightarrow$ 4;Q $\rightarrow$ 3;R $\rightarrow$ 2;S $\rightarrow$ 1}
\end{qoptions}
\qn{12} Match the following List-I and List-II :\par
{\footnotesize\setlength{\tabcolsep}{1.5pt}\noindent\begin{tabular}{|c|L{6.0cm}|c|L{0.8cm}|}
\hline
 & \multicolumn{1}{>{\centering\arraybackslash}p{6.0cm}|}{List-I} & & \multicolumn{1}{>{\centering\arraybackslash}p{0.8cm}|}{List-II} \\ \hline
(P) & If A and B are non-singular real matrices such that det(A) = 3 and A$^{-1}$ + B$^{-1}$ = (A + B)$^{-1}$, then det(B) = & (1) & 190 \\ \hline
(Q) & Let A $=\begin{pmatrix}a&x&p\\y&q&b\\r&c&z\end{pmatrix}$ and B $=\begin{pmatrix}0&0&1\\0&1&0\\1&0&0\end{pmatrix}$ (a,b,c,x,y,z,p,q,r $\in$ N). If tr(AB + AB$^3$ + AB$^5$+....+ AB$^{19}$) = 210, then the number of ordered triplets (p,q,r) is ........(tr(A) denotes sum of principal diagonal elements of A) & (2) & 3 \\ \hline
(R) & Consider A $=\begin{pmatrix}1&2&4\\0&1&3\\0&0&1\end{pmatrix}$ B $=\begin{pmatrix}2&4&5\\4&0&0\\5&0&0\end{pmatrix}$. If C = [c$_{ij}$]$_{3\times3}$ = A$^{20}$ + B and c$_{22}$c$_{33}$ $-$ c$_{23}$c$_{32}$ = 2$^m$, then m = & (3) & 5 \\ \hline
(S) & Consider three matrices A = [a$_{ij}$]$_{3\times5}$, B = [b$_{ij}$]$_{5\times9}$, and C = A $\times$ B = [c$_{ij}$]$_{3\times9}$. If is given that a$_{ij}$ = $\dfrac i{3j-2}$ and b$_{ij}$ = $\dfrac j{3i+1}$, then value of the element c$_{28}$ is & (4) & 0 \\ \hline
\end{tabular}}\par
\begin{qoptions}
\opt{A}{P $\rightarrow$ 2;Q $\rightarrow$ 1;R $\rightarrow$ 4;S $\rightarrow$ 3}
\opt{B}{P $\rightarrow$ 4;Q $\rightarrow$ 3;R $\rightarrow$ 1;S $\rightarrow$ 2}
\opt{C}{P $\rightarrow$ 3;Q $\rightarrow$ 4;R $\rightarrow$ 2;S $\rightarrow$ 1}
\opt{D}{P $\rightarrow$ 4;Q $\rightarrow$ 3;R $\rightarrow$ 2;S $\rightarrow$ 1}
\end{qoptions}
\sect{SECTION-II}
\qn{1} Suppose A $\in$ \{a$_1$,a$_2$,a$_3$,a$_4$,a$_5$\} where a$_1$ $<$ a$_2$ $<$ a$_3$ $<$ a$_4$ $<$ a$_5$ and B $\in$ \{1,2,3,4,....k\} then, match List-I with List-II\par
{\footnotesize\setlength{\tabcolsep}{1.5pt}\noindent\begin{tabular}{|c|L{5.2cm}|c|L{1.0cm}|}
\hline
 & \multicolumn{1}{>{\centering\arraybackslash}p{5.2cm}|}{List-I} & & \multicolumn{1}{>{\centering\arraybackslash}p{1.0cm}|}{List-II} \\ \hline
(P) & No. of injective functions in f : A $\rightarrow$ B such that f(a$_i$) $\ne$ i for k = 6 & (1) & 2002 \\ \hline
(Q) & No. of non-decreasing function in f : A $\rightarrow$ B for k = 11 & (2) & 309 \\ \hline
(R) & No. of strictly increasing functions in f : A $\rightarrow$ B such that f(a$_i$) $\ne$ i for k = 9 & (3) & 3003 \\ \hline
(S) & No. of non-increasing function in f : A $\rightarrow$ B for k = 10 & (4) & 56 \\ \hline
\end{tabular}}\par
\qn{2} \par
{\footnotesize\setlength{\tabcolsep}{1.5pt}\noindent\begin{tabular}{|c|L{5.0cm}|c|L{2.3cm}|}
\hline
 & \multicolumn{1}{>{\centering\arraybackslash}p{5.0cm}|}{List-I} & & \multicolumn{1}{>{\centering\arraybackslash}p{2.3cm}|}{List-II} \\ \hline
(P) & \resizebox{5.0cm}{!}{$\dfrac1n\left(\displaystyle\sum_{0\le i<j\le n}\sum\left((i+j)\left({}^n\mathrm C_i\,{}^n\mathrm C_j\right)\right)\right)$} is equal to & (1) & n.(${}^{2n}$C$_n$) \\ \hline
(Q) & \resizebox{4.6cm}{!}{$\left(\displaystyle\sum_{0\le i<j\le n}\sum\left(\left({}^n\mathrm C_i\right)^2+\left({}^n\mathrm C_j\right)^2\right)\right)$} is equal to & (2) & $\dfrac12\left(2^{2n}-{}^{2n}\mathrm C_n\right)$ \\ \hline
(R) & \resizebox{5.0cm}{!}{$\displaystyle\sum_{0\le i<j\le n}\sum(i+j)\left({}^n\mathrm C_i-{}^n\mathrm C_j\right)$} is equal to & (3) & [(n + 1) ${}^{2n}$C$_n$ $-$ 2$^{2n}$] n \\ \hline
(S) & $\displaystyle\sum_{0\le i<j\le n}\sum{}^n\mathrm C_i\,{}^n\mathrm C_j$ is equal to & (4) & $\displaystyle\sum_{r=0}^{n-1}{}^{2n}\mathrm C_r$ \\ \hline
\end{tabular}}\par
\qn{3} There are 16 equally skilled players S$_1$,S$_2$,S$_3$,......,S$_{16}$ playing a knockout tournament. They are divided into 8 pairs at random. From each pair a winner is decided on the basis of a game played between the two players of the pair, and they enter into the next round the tournament proceed in the similar way.\par
{\footnotesize\setlength{\tabcolsep}{1.5pt}\noindent\begin{tabular}{|c|L{5.6cm}|c|L{0.6cm}|}
\hline
 & \multicolumn{1}{>{\centering\arraybackslash}p{5.6cm}|}{List-I} & & \multicolumn{1}{>{\centering\arraybackslash}p{0.6cm}|}{List-II} \\ \hline
(P) & Probability that exactly one of S$_1$ and S$_2$ is among the 4 winners of 2$^{\mathrm{nd}}$ round is & (1) & $\dfrac1{30}$ \\ \hline
(Q) & The probability that S$_1$ wins the tournament given that S$_2$ enters in the semifinals is & (2) & $\dfrac1{20}$ \\ \hline
(R) & The probability that S$_1$ wins the tournament given that S$_2$ enters into final is & (3) & $\dfrac7{30}$ \\ \hline
(S) & The probability that S$_1$ is among the 4 winners of 2nd round given that S$_2$ wins the first round is & (4) & $\dfrac25$ \\ \hline
\end{tabular}}\par
\qn{4} \par
{\footnotesize\setlength{\tabcolsep}{1.5pt}\noindent\begin{tabular}{|c|L{5.6cm}|c|L{1.0cm}|}
\hline
 & \multicolumn{1}{>{\centering\arraybackslash}p{5.6cm}|}{List-I} & & \multicolumn{1}{>{\centering\arraybackslash}p{1.0cm}|}{List-II} \\ \hline
(P) & The least value of 8$|$z $-$ 7$|$ + 6$|$z $-$ 5$|$ ; z $\in$ C is & (1) & 12 \\ \hline
(Q) & The least value of $\left|\dfrac{z-7+11i}{\sqrt2}\right|+\left|\dfrac{z+5-i}{\sqrt2}\right|+\left|\dfrac{z+4}{\sqrt2}\right|$ ; z $\in$ C is & (2) & 13 \\ \hline
(R) & The the least value of $|$z $-$ 3 $-$ 4i$|$ + $|$z $-$ i$|$ + $|$z$|$ + $|$z $-$ 1$|$ ; z $\in$ C is & (3) & 10 \\ \hline
(S) & The least value of $|$z $-$ 3$|^2$ + $|$z $-$ 5 + 2i$|^2$ + $|$z $-$ 1 + i$|^2$ is & (4) & 5 + $\sqrt2$ \\ \hline
\end{tabular}}\par
\qn{5} A variable plane cuts the x-axis, y-axis and z-axis at the points A,B and C respectively such that the volume of the tetrahedron OABC remains constant equal 32 cubic unit and O is the origin of the coordinate system.\par
{\footnotesize\setlength{\tabcolsep}{1.5pt}\noindent\begin{tabular}{|c|L{4.0cm}|c|L{3.3cm}|}
\hline
 & \multicolumn{1}{>{\centering\arraybackslash}p{4.0cm}|}{List-I} & & \multicolumn{1}{>{\centering\arraybackslash}p{3.3cm}|}{List-II} \\ \hline
(P) & The locus of the centroid of the tetrahedron is & (1) & xyz = 24 \\ \hline
(Q) & The locus of the point equidistant from O,A,B and C is & (2) & (x$^2$ + y$^2$ + z$^2$)$^3$ = 192xyz \\ \hline
(R) & The locus of the foot of perpendicular from origin to the plane is & (3) & xyz = 3 \\ \hline
(S) & If PA, PB and PC are mutually perpendicular then the locus of P is & (4) & (x$^2$ + y$^2$ + z$^2$)$^3$ = 1536xyz \\ \hline
\end{tabular}}\par
\qn{6} Let P(x) = x$^5$ + x$^2$ + 1 has zeros $\alpha_1,\alpha_2,\alpha_3,\alpha_4,\alpha_5$ and Q(x) = x$^2$ $-$ 2, then match List-I with List-II :\par
{\footnotesize\setlength{\tabcolsep}{1.5pt}\noindent\begin{tabular}{|c|L{4.6cm}|c|L{1.0cm}|}
\hline
 & \multicolumn{1}{>{\centering\arraybackslash}p{4.6cm}|}{List-I} & & \multicolumn{1}{>{\centering\arraybackslash}p{1.0cm}|}{List-II} \\ \hline
(P) & $\displaystyle\prod_{i=1}^5\mathrm Q(\alpha_i)=$ & (1) & 20 \\ \hline
(Q) & $\displaystyle\sum_{i=1}^5\mathrm Q(\alpha_i)=$ & (2) & $-$23 \\ \hline
(R) & $\displaystyle\sum_{1\le i<j\le5}\mathrm Q(\alpha_i)\mathrm Q(\alpha_j)=$ & (3) & 40 \\ \hline
(S) & $\displaystyle\sum_{i=1}^5\mathrm Q^2(\alpha_i)=$ & (4) & $-$10 \\ \hline
\end{tabular}}\par
\end{multicols}
\end{document}
